% Blueprint note: Fixed upstream source versions: \cite{KIPSource5aa845241533,KIPSourcebbac1fc05843}. Interface scope notes are retained as native hgraph node comments; historical migration records are local archives.

% Curated mathematical authoring. Lean signatures remain the scope authority.

\section{Classical input and realization}

\begin{proposition}[Low-weight sphere comparison isomorphisms]
\label{sem:Literature:low-sphere-comparisons}
\level{coarse}
\uses{sem:AdamsSynthetic:nu,paper:3.2}
(Project construction from the unit and suspension comparisons of $\nu$;
cf.\ \cite{LWX126}, Definition 3.2.)\par{}

For $k=1,2,3$, the fixed unit and suspension comparisons give isomorphisms
\[\nu(S^k)\cong S^{k,k}.\]
These are the three specified comparisons, rather than an assertion for
every integral weight. The sphere objects themselves are defined in
\ref{paper:3.2}.
\lean{KIP126.Literature.nuSphereOne,KIP126.Literature.nuSphereTwo,KIP126.Literature.nuSphereThree}
\end{proposition}

\begin{theorem}[Connectivity and mod-two homology]
\label{sem:Literature:connected-homology}
\level{coarse}
\uses{sem:AdamsSynthetic:categories}
(Adapted from \cite{SchwedeSymmetricSpectra}, Prop. II.5.22; \cite{EKMM}, Thm. IV.4.5.)\par{}

The sphere is $(-1)$-connected. If $X$ is $n$-connected, then $H_i(X;\mathbb F_2)=0$ for $i\le n$. There is an additive universal-coefficient equivalence
\[H^n(X;\mathbb F_2)\cong\operatorname{Hom}(H_{-n}(X;\mathbb F_2),\mathbb F_2).\]
Both groups use the same classical category and coefficient object.
\lean{KIP126.Literature.Standard.sphere_connected,KIP126.Literature.Standard.connected_homology_vanishing,KIP126.Literature.Standard.universal_coefficient}
% Blueprint note: Source verification (2026-10-18): Source adaptation checked. The notation H^n here means the project legacy group [Sigma^n X,HF2], i.e. ordinary H^{-n}. The connectivity definition and this sign convention are project definitions. The UCF uses the full dual and requires no finite-type assumption. This verifies the mathematical source and the stated adaptation, not the pending comparison with the fixed project model. No Lean proof obligation is discharged by this audit. Primary references: [SchwedeSymmetricSpectra](https://www.math.uni-bonn.de/~schwede/SymSpec-v3.pdf): Example I.1.8; proof of Proposition II.5.21; pages [12, 264] [SchwedeSymmetricSpectra](https://www.math.uni-bonn.de/~schwede/SymSpec-v3.pdf): Proposition II.5.22; Example I.1.14; Definition II.6.21; pages [16, 17, 264, 265, 278, 279] [EKMM](https://www.math.uchicago.edu/~may/BOOKS/EKMM.pdf): Proposition IV.1.2(i), after shifting X by -(n+1); pages [72] [SchwedeSymmetricSpectra](https://www.math.uni-bonn.de/~schwede/SymSpec-v3.pdf): Proposition II.6.28 and Equation (6.27); pages [281] [EKMM](https://www.math.uchicago.edu/~may/BOOKS/EKMM.pdf): Theorem IV.4.5, specialized R=N=HF2; pages [82]
% Blueprint note: Detailed source attribution: (Adapted from \cite{SchwedeSymmetricSpectra}, Example I.1.8, Proposition II.5.22, and Definition II.6.21; and from \cite{EKMM}, Theorem IV.4.5, specialized to $H\mathbb F_2$. The cohomological degree uses the project convention $H^n_{\mathrm{legacy}}(X)=H^{-n}(X;\mathbb F_2)$.)
\end{theorem}

\begin{definition}[Realization in sphere coordinates]
\label{sem:Literature:realization-coordinates}
\level{coarse}
\uses{sem:AdamsSynthetic:recovery,paper:3.2,paper:3.3}
(Project construction; cf.\ \cite{Pst}, Prop. 4.40.)\par{}

Let $\operatorname{Re}$ be the additive realization functor with sphere isomorphisms $S^m\cong\operatorname{Re}(S^{m,w})$. Its isomorphism at the unit is the $\nu$--realization comparison. In these coordinates, realization sends each sphere map $\lambda^k$ to the identity. Write $\operatorname{Re}(a)$ for the induced realization of a sphere class or a class on $\nu X$.
\lean{KIP126.Literature.Standard.realization_additive,KIP126.Literature.Standard.realization_sphere,KIP126.Literature.Standard.realization_unit,KIP126.Literature.Standard.realization_lambda,KIP126.Literature.Standard.realizationData,KIP126.Literature.Standard.realizeSphere,KIP126.Literature.Standard.realizeNu}
% Blueprint note: Source verification (2026-10-18): Source adaptation checked. The canonical synthetic realization has these coordinates; equalities for the project-selected sphere, unit, shift, and lambda comparisons remain explicit source-to-model obligations. These are project definitions packaging coordinates and transporting homotopy classes; the article does not contain these names or choices. Additivity is proved in Lean from the actual realization adjunction and biproduct preservation, rather than quoted as a separate numbered theorem. Primary references: [Pst](https://arxiv.org/pdf/1803.01804v3): Definition 4.6; Proposition 4.40; Remark 4.41; pages 39, 49; arXiv:1803.01804v3, 10 November 2022
% Blueprint note: Detailed source attribution: (Adapted from \cite{Pst}, Definition 4.6, Proposition 4.40, and Remark 4.41.)
\end{definition}

\begin{proposition}[The Hopf classes and their synthetic normalization]
\label{sem:Literature:hopf-classes}
\level{coarse}
\uses{sem:AdamsSynthetic:standard-classes,sem:AdamsSynthetic:normalized-map,sem:Literature:realization-coordinates,sem:AdamsSynthetic:e2-comparison}
(Adapted from \cite{Adams58}, Lem. 2.3; \cite{LWX126}, Prop. 3.16.)\par{}

The classical maps $\eta$ and $\nu$ are detected by the nonzero permanent cycles $h_1\in E_2^{1,2}$ and $h_2\in E_2^{1,4}$, respectively. The normalized synthetic lift $\eta\in\pi_{1,2}S^{0,0}$ has filtration exponent one, first-quotient label $h_1$, and realization equal to the classical $\eta$.
\lean{KIP126.Literature.Standard.eta,KIP126.Literature.Standard.eta_normalized,KIP126.Literature.Standard.eta_detection,KIP126.Literature.Standard.eta_exponent,KIP126.Literature.Standard.eta_permanent,KIP126.Literature.Standard.eta_label,KIP126.Literature.Standard.eta_realization,KIP126.Literature.Standard.nu_detection,KIP126.Literature.Standard.nu_permanent}
% Blueprint note: Source verification (2026-10-18): Source adaptation checked. The classical detections and nonzero permanence are ordinary Hopf facts. The normalized synthetic map is the project chosen lift with exponent e(eta)=1, not a literal named theorem in Adams60. Realization uses the specified nu-unit and lambda-inverting comparison. This verifies the mathematical source and the stated adaptation, not the pending comparison with the fixed project model. No Lean proof obligation is discharged by this audit. Primary references: [Adams58](https://www.sas.rochester.edu/mth/sites/doug-ravenel/otherpapers/Adams-SS.pdf): Lemma 2.3 and the following m=0 paragraph; pages 184 [Adams60](https://www.maths.ed.ac.uk/~v1ranick/papers/adams1.pdf): Section 1.1, after Theorem 1.1.1; pages 20 [LWX126](https://arxiv.org/abs/2412.10879): Proposition 3.16; Notation 3.19; Notation 7.2; pages [18, 19, 42] [Pst](https://arxiv.org/pdf/1803.01804v3): Proposition 4.40 and Remark 4.41; pages 49
% Blueprint note: Detailed source attribution: (Adapted from \cite{Adams58}, Lemma 2.3, p.~184; \cite{Adams60}, Section 1.1, p.~20; and \cite{LWX126}, Proposition 3.16 and Notation 3.19.)
\end{proposition}

\begin{proposition}[A detected Kervaire class in stem 62]
\label{sem:Literature:theta-five}
\level{coarse}
\uses{sem:AdamsSynthetic:standard-classes,sem:AdamsSynthetic:selected-convergence,sem:Literature:realization-coordinates,sem:AdamsSynthetic:e2-comparison}
(Adapted from \cite{BJMtheta5}, Thm. 2.1; \cite{Xu}, Cor. 1.3.)\par{}

There is a classical class $\theta_5\in\pi_{62}S^0$ of order two detected by the nonzero permanent cycle $h_5^2\in E_2^{2,64}$. Fix such a class. There is a synthetic lift $a\in\pi_{62,64}S^{0,0}$ with first-quotient label $h_5^2$ and $\operatorname{Re}(a)=\theta_5$.
% Blueprint note: The existential synthetic detection statement does not identify its witness with any separately chosen synthetic theta-five.
\lean{KIP126.Literature.Standard.theta5,KIP126.Literature.Standard.h5Square_permanent,KIP126.Literature.Standard.theta5_detection,KIP126.Literature.Standard.theta5_order_two,KIP126.Literature.Standard.theta5_lift,KIP126.Literature.Standard.theta5_detected_exists}
% Blueprint note: Source verification (2026-10-18): Source adaptation checked. Select the order-two 2-primary theta5 guaranteed by Xu, rather than asserting that every ordinary integral representative detected mod two has order two. Bidegrees (s,t)=(2,64) give synthetic (62,64). The existence of a prescribed-realization lift is derived using the project comparison theorem. This verifies the mathematical source and the stated adaptation, not the pending comparison with the fixed project model. No Lean proof obligation is discharged by this audit. Primary references: [BJMtheta5](https://www.sas.rochester.edu/mth/sites/doug-ravenel/otherpapers/barjoma.pdf): Theorem 2.1 and its immediate conclusion; pages [536, 537] [Xu](https://msp.org/gt/2016/20-3/gt-v20-n3-p08-p.pdf): Corollary 1.3; pages 1612 [IWX](https://arxiv.org/pdf/2001.04511v3): Corollary 1.3, Table 1, row k=62; pages [4, 6] [LWX126](https://arxiv.org/abs/2412.10879): Theorem 3.8(1), Proposition 3.11, and Notation 7.2; pages [16, 17, 42]
% Blueprint note: Detailed source attribution: (Adapted from \cite{BJMtheta5}, Theorem 2.1, pp.~536--537; \cite{Xu}, Corollary 1.3, p.~1612; and \cite{LWX126}, Theorem 3.8 and Notation 7.2.)
\end{proposition}

\begin{proposition}[The 2-primary stem and the filtration of choices]
\label{sem:Literature:theta-indeterminacy}
\level{coarse}
\uses{sem:Literature:theta-five,sem:AdamsSynthetic:selected-convergence}
(Adapted from \cite{IWX}, Table 1; \cite{LWX126}, proof of Lem. 7.10.)\par{}

Every class $a\in\pi_{62}S^0$ annihilated by a power of two satisfies $2a=0$. If $a,b\in\pi_{62}S^0$ are detected by $h_5^2$, then $a-b$ has Adams filtration at least six.
\lean{KIP126.Literature.Standard.stem62_exponent_two,KIP126.Literature.Standard.theta5_filtration_gap}
% Blueprint note: Source verification (2026-10-18): Source adaptation checked. Only the 2-primary subgroup has exponent two. The ordinary 62-stem also has an odd-primary summand. The filtration assertion concerns two representatives with the same h5^2 detection; their difference may include an infinitely filtered odd-primary component. This verifies the mathematical source and the stated adaptation, not the pending comparison with the fixed project model. No Lean proof obligation is discharged by this audit. Primary references: [IWX](https://arxiv.org/pdf/2001.04511v3): Corollary 1.3, Table 1, row k=62; pages [4, 6] [LWX126](https://arxiv.org/abs/2412.10879): Proof of Lemma 7.10; pages 45 [IWXCharts](https://zenodo.org/records/6987157/files/Adams-classical-Einfty.csv?download=1): Adams-classical-Einfty.csv, physical rows 204--207 (including header)
% Blueprint note: Detailed source attribution: (Adapted from \cite{IWX}, Corollary 1.3 and Table 1, row 62; and \cite{LWX126}, proof of Lemma 7.10, p.~45.)
\end{proposition}

\begin{theorem}[The positive-stem Adams vanishing range]
\label{sem:Literature:sphere-vanishing}
\level{coarse}
\uses{sem:AdamsSynthetic:families}
(Adapted from \cite{MathewTorsion}, Prop. 3.3(a).)\par{}

In the classical sphere Adams spectral sequence,
\[0<t-s<2s-3\quad\Longrightarrow\quad E_2^{s,t}(S^0)=0.\]
\lean{KIP126.Literature.Standard.sphereVanishing}
% Blueprint note: Source verification (2026-10-18): Source adaptation checked. This is the classical mod-two sphere E2-term, after the existing project identification with Steenrod Ext. The displayed assumptions imply s>=3, so the source hypothesis 0<s is automatic. This verifies the mathematical source and the stated adaptation, not the pending comparison with the fixed project model. No Lean proof obligation is discharged by this audit. Primary references: [MathewTorsion](https://msp.org/agt/2016/16-2/agt-v16-n2-p12-s.pdf): Proposition 3.3(a); pages 1031
% Blueprint note: Detailed source attribution: (Adapted from \cite{MathewTorsion}, Proposition 3.3(a), p.~1031.)
\end{theorem}

\begin{theorem}[The Hopf-invariant-one differentials]
\label{sem:Literature:one-line-differentials}
\level{coarse}
\uses{sem:AdamsSynthetic:standard-classes,sem:AdamsSynthetic:selected-convergence}
(Adapted from \cite{Adams58}, Lem. 2.3; \cite{LWX126}, Introduction.)\par{}

The class $h_0$ detects the classical degree-two map. For each integer $j\ge4$, there is a nonzero differential
\[d_2(h_j)=h_0h_{j-1}^2\in E_2^{3,1+2^j}(S^0).\]
\lean{KIP126.Literature.Standard.two_detection,KIP126.Literature.Standard.adamsOneLine_d2}
% Blueprint note: Source verification (2026-10-18): Source adaptation checked. The target is the actual product h0 h_(j-1)^2, not a separate arbitrary class. For j>=4 the Nat expression j-1+1 equals j, so the Lean target degree is (3,1+2^j). This verifies the mathematical source and the stated adaptation, not the pending comparison with the fixed project model. No Lean proof obligation is discharged by this audit. Primary references: [Adams58](https://www.sas.rochester.edu/mth/sites/doug-ravenel/otherpapers/Adams-SS.pdf): Lemma 2.3 and the following m=0 paragraph; pages 184 [LWX126](https://arxiv.org/abs/2412.10879): Introduction, unnumbered displayed Hopf-invariant-one differential; pages 2
% Blueprint note: Detailed source attribution: (Adapted from \cite{Adams58}, Lemma 2.3, p.~184; and \cite{LWX126}, Introduction, p.~2, the displayed Hopf-invariant-one differential.)
\end{theorem}

\section{The tmf detector}

\begin{definition}[The chosen tmf object, algebra, and unit]
\label{sem:Literature:tmf-unit}
\level{coarse}
\uses{sem:AdamsSynthetic:categories,frontier:AdamsSynthetic:0895d4999173}
(Project construction; cf.\ \cite{BR21}, §0.1.)\par{}

Let $\mathit{tmf}$ be the classical spectrum of topological modular forms, with its algebra structure and unit $u:S^0\to\mathit{tmf}$. Fix an isomorphism $\mathit{tmf}\cong D$ carrying $u$ to the unit of the detector $D$.
\lean{KIP126.Literature.Standard.tmf,KIP126.Literature.Standard.tmfIso,KIP126.Literature.Standard.tmf_algebra,KIP126.Literature.Standard.tmfUnit,KIP126.Literature.Standard.tmfUnit_agrees,KIP126.Literature.Standard.tmf_algebra_unit}
% Blueprint note: Source verification (2026-10-18): Partial source correspondence. The cited papers concern the usual classical spectrum or its 2-primary/2-completed version. Recognition of the fixed project category, Adams tower, coefficient object, and selected Lean constants is a separate, currently pending formal comparison; this audit does not prove it. Primary references: [BR21](https://www.sas.rochester.edu/mth/sites/doug-ravenel/otherpapers/rognes-bruner-tmf-book.pdf): Section0.1,p1 (source ring unit only); Section0.1,pp1–2 (source object only); Section0.1,pp1–2; Section0.1,p1;Section1.2,p64; Section0.1,p1 (underlying ring spectrum); Author preliminary version of AMS Mathematical Surveys and Monographs 253 (2021), downloaded 2026-10-18; printed page numbers cited, PDF index = printed page + 17 in the body.
% Blueprint note: Detailed source attribution: (Adapted from \cite{BR21}, Sections 0.1 and 1.2; the detector and unit comparisons are project data.)
\end{definition}

\begin{proposition}[Connectivity, finite type, and convergence for tmf]
\label{sem:Literature:tmf-convergence}
\level{coarse}
\uses{sem:Literature:tmf-unit,sem:AdamsSynthetic:families,sem:AdamsSynthetic:selected-convergence}
(Adapted from \cite{BR21}, §1.2.)\par{}

The spectrum $\mathit{tmf}$ is connective and has finite mod-two homology in each degree. Its classical Adams tower converges to its homotopy groups with their Adams filtration. Transport along $\mathit{tmf}\cong D$ preserves tower lifts and associated-graded representatives.
\lean{KIP126.Literature.Standard.tmf_connective,KIP126.Literature.Standard.tmf_finiteModTwoType,KIP126.Literature.Standard.tmf_convergence}
% Blueprint note: Source verification (2026-10-18): Partial source correspondence. The cited papers concern the usual classical spectrum or its 2-primary/2-completed version. Recognition of the fixed project category, Adams tower, coefficient object, and selected Lean constants is a separate, currently pending formal comparison; this audit does not prove it. Primary references: [BR21](https://www.sas.rochester.edu/mth/sites/doug-ravenel/otherpapers/rognes-bruner-tmf-book.pdf): Section0.3,p5;Section0.5,pp8–9;Section1.2,p64; Section0.1,p1;Section0.5,pp8–9;Section1.2,p64; Author preliminary version of AMS Mathematical Surveys and Monographs 253 (2021), downloaded 2026-10-18; printed page numbers cited, PDF index = printed page + 17 in the body.
% Blueprint note: Detailed source attribution: (Adapted from \cite{BR21}, Section 0.3, p. 5, and Section 1.2, p. 64; the completed-model comparison is a project obligation.)
\end{proposition}

\begin{proposition}[The class g and the detection of bar-kappa]
\label{sem:Literature:kappa-detection}
\level{coarse}
\uses{sem:AdamsSynthetic:families,sem:AdamsSynthetic:selected-convergence}
(Adapted from \cite{BR21}, Thm. 11.61(20).)\par{}

The group $E_2^{4,24}(S^0)$ has a unique nonzero element $g$. It detects a class $\bar\kappa\in\pi_{20}S^0$, which we fix.
\lean{KIP126.Literature.Standard.g,KIP126.Literature.Standard.g_nonzero,KIP126.Literature.Standard.g_unique,KIP126.Literature.Standard.kappaBar,KIP126.Literature.Standard.kappaBar_detection}
% Blueprint note: Source verification (2026-10-18): Source adaptation checked. The cited papers concern the usual classical spectrum or its 2-primary/2-completed version. Recognition of the fixed project category, Adams tower, coefficient object, and selected Lean constants is a separate, currently pending formal comparison; this audit does not prove it. Primary references: [BR21](https://www.sas.rochester.edu/mth/sites/doug-ravenel/otherpapers/rognes-bruner-tmf-book.pdf): Proposition1.4;Table1.1,p54, row(stem20,filtration4), cocycle4_8;Figure1.9,p56; Theorem11.61,case(20),p467;Table11.3,p458; Proposition1.4;Table1.1,p54;Remark1.5,pp56–57;Figure1.9,p56; Proposition1.4;Table1.1,p54;Figure1.9,p56; Author preliminary version of AMS Mathematical Surveys and Monographs 253 (2021), downloaded 2026-10-18; printed page numbers cited, PDF index = printed page + 17 in the body.
% Blueprint note: Detailed source attribution: (Adapted from \cite{BR21}, Proposition 1.4, Table 1.1, and Theorem 11.61(20).)
\end{proposition}

\begin{proposition}[The class Delta h-one g and the detection of w]
\label{sem:Literature:w-detection}
\level{coarse}
\uses{sem:AdamsSynthetic:families,sem:AdamsSynthetic:selected-convergence}
(Adapted from \cite{BR21}, Thm. 11.59.)\par{}

The group $E_2^{9,54}(S^0)$ has a unique nonzero element $\Delta h_1g$. It detects a class $w\in\pi_{45}S^0$, which we fix.
\lean{KIP126.Literature.Standard.deltaH1g,KIP126.Literature.Standard.deltaH1g_nonzero,KIP126.Literature.Standard.deltaH1g_unique,KIP126.Literature.Standard.wClass,KIP126.Literature.Standard.w_detection}
% Blueprint note: Source verification (2026-10-18): Source adaptation checked. The cited papers concern the usual classical spectrum or its 2-primary/2-completed version. Recognition of the fixed project category, Adams tower, coefficient object, and selected Lean constants is a separate, currently pending formal comparison; this audit does not prove it. Primary references: [BR21](https://www.sas.rochester.edu/mth/sites/doug-ravenel/otherpapers/rognes-bruner-tmf-book.pdf): BR21 Proposition1.4;Table1.1,p56,row(stem45,filtration9), classical Ext namew,cocycle9_20;Remark1.5,pp56–57;Figure1.10,p57; IWX Section1.5,pp117–118;Table4,p215,row(45,9,24),τΔh1g,alternate namew; BR21 Proposition1.4;Table1.1,p56;Remark1.5,pp56–57;Figure1.10,p57; IWX Table4,p215,row(45,9,24); BR21 Theorem11.59,pp457–458;Table11.3,p458,roww in(stem45,filtration9); BMQ Section1,pp3–4,names following Figure1.1; BR21 Proposition1.4;Table1.1,p56;Figure1.10,p57; BR21 Theorem11.59,pp457–458;Table11.3,p458,roww; Author preliminary version of AMS Mathematical Surveys and Monographs 253 (2021), downloaded 2026-10-18; printed page numbers cited, PDF index = printed page + 17 in the body. [IWX](https://www.numdam.org/item/10.1007/s10240-023-00139-1.pdf): BR21 Proposition1.4;Table1.1,p56,row(stem45,filtration9), classical Ext namew,cocycle9_20;Remark1.5,pp56–57;Figure1.10,p57; IWX Section1.5,pp117–118;Table4,p215,row(45,9,24),τΔh1g,alternate namew; BR21 Proposition1.4;Table1.1,p56;Remark1.5,pp56–57;Figure1.10,p57; IWX Table4,p215,row(45,9,24); BR21 Theorem11.59,pp457–458;Table11.3,p458,roww in(stem45,filtration9); BMQ Section1,pp3–4,names following Figure1.1; BR21 Proposition1.4;Table1.1,p56;Figure1.10,p57; BR21 Theorem11.59,pp457–458;Table11.3,p458,roww; Published IHES137 (2023),107–243; PDF index=printed page−106. [tmf](https://arxiv.org/pdf/2011.08956v4): BR21 Theorem11.59,pp457–458;Table11.3,p458,roww in(stem45,filtration9); IWX Table4,p215,row(45,9,24); BMQ Section1,pp3–4,names following Figure1.1; arXiv:2011.08956v4, 20 November 2021; published Geom. Topol.27 (2023)2763–2831. Locators below use v4 page numbering.
% Blueprint note: Detailed source attribution: (Adapted from \cite{BR21}, Proposition 1.4, Table 1.1, Theorem 11.59, and Table 11.3; notation as in \cite{IWX}, Table 4.)
\end{proposition}

\begin{proposition}[The tmf Hurewicz input in stems 62 and 125]
\label{sem:Literature:tmf-hurewicz}
\level{coarse}
\uses{sem:Literature:tmf-unit,sem:Literature:tmf-convergence,sem:Literature:kappa-detection,sem:Literature:w-detection}
(Adapted from \cite{tmf}, Fig. 1.1 and Thm. 1.2.)\par{}

One has $\pi_{62}\mathit{tmf}=0$ and $E_2^{s,63+s}(\mathit{tmf})=0$ for $s\le0$. Let $v=\bar\kappa^4w\in\pi_{125}S^0$. For every positive odd integer $n$,
\[n\,u_*(v)\ne0\quad\text{in }\pi_{125}\mathit{tmf}.\]
\lean{KIP126.Literature.Standard.tmf_vanishing62,KIP126.Literature.Standard.tmf_low_filtration63,KIP126.Literature.Standard.high125,KIP126.Literature.Standard.tmf_high125_nonzero}
% Blueprint note: Source verification (2026-10-18): Source adaptation checked. The cited papers concern the usual classical spectrum or its 2-primary/2-completed version. Recognition of the fixed project category, Adams tower, coefficient object, and selected Lean constants is a separate, currently pending formal comparison; this audit does not prove it. Primary references: [BR21](https://www.sas.rochester.edu/mth/sites/doug-ravenel/otherpapers/rognes-bruner-tmf-book.pdf): Section1.2,p64,E2=Ext_A(2)(F2,F2);Section0.3,p5;Section0.5,pp8–9; BMQ Figure1.1,p3,nonzeroκbar^4w atstem125;Section1,pp3–4;Theorem1.2,p4;discussionpp5–6; BR21 Theorem9.27,pp338–340;Table9.4,stem125row; Author preliminary version of AMS Mathematical Surveys and Monographs 253 (2021), downloaded 2026-10-18; printed page numbers cited, PDF index = printed page + 17 in the body. [tmf](https://arxiv.org/pdf/2011.08956v4): BMQ Figure1.1,p3,nonzeroκbar^4w atstem125;Section1,pp3–4;Theorem1.2,p4;discussionpp5–6; BR21 Theorem9.27,pp338–340;Table9.4,stem125row; BMQ Figure1.1,p3;Section1,pp3–4; Figure1.1,p3;chart conventions in Section1,pp2–4; arXiv:2011.08956v4, 20 November 2021; published Geom. Topol.27 (2023)2763–2831. Locators below use v4 page numbering.
% Blueprint note: Detailed source attribution: (Adapted from \cite{tmf}, Figure 1.1 and Theorem 1.2, and \cite{BR21}, Section 1.2; the odd-multiple claim is a 2-local corollary.)
\end{proposition}

\begin{theorem}[The power-operation differential and detector transport]
\label{sem:Literature:tmf-d-three}
\level{coarse}
\uses{sem:Literature:tmf-unit,sem:AdamsSynthetic:families}
(Adapted from \cite{BR21}, Thm. 5.18.)\par{}

For the tmf classes $v_2^{16}\in E_2^{16,112}$ and $\beta^5g\in E_2^{19,114}$,
\[d_3(v_2^{16})=\beta^5g.\]
Under $\mathit{tmf}\cong D$, the images of these classes satisfy the same differential.
% Blueprint note: BR21, Theorem 5.18 and Table 5.4, writes the target as $\beta g^4$; Table 3.5 gives $\beta^4=g^3$, so this equals the AimPaper expression $\beta^5g$. Identifying the specified formal target with this source label remains an external-input obligation.
\lean{KIP126.Literature.Standard.tmf_br21,KIP126.Literature.Standard.detector_d3}
% Blueprint note: Source verification (2026-10-18): Source adaptation checked. The cited papers concern the usual classical spectrum or its 2-primary/2-completed version. Recognition of the fixed project category, Adams tower, coefficient object, and selected Lean constants is a separate, currently pending formal comparison; this audit does not prove it. Primary references: [BR21](https://www.sas.rochester.edu/mth/sites/doug-ravenel/otherpapers/rognes-bruner-tmf-book.pdf): BR21 Theorem5.18;Table5.4,p196;Remark5.9,p190;Table3.5,p151;Theorem1.20 proof,p77; AimPaper Appendix,paragraph immediately beforeTable5; Theorem5.18;Table5.4,p196;Proposition5.8(4)andRemark5.9,p190;Table3.5,p151;Table1.3,p65; Source differential:Theorem5.18andTable5.4,p196;projectderivation:Standard.lean:510–514; BR21 Proposition1.12;Table1.3,p65,restrictionw2↦v2^8;Table5.4,p196,roww2²; AimPaper Appendix,paragraph immediately beforeTable5: d3(v2^16)=β^5g; Author preliminary version of AMS Mathematical Surveys and Monographs 253 (2021), downloaded 2026-10-18; printed page numbers cited, PDF index = printed page + 17 in the body. [LWX126](https://arxiv.org/abs/2412.10879): BR21 Theorem5.18;Table5.4,p196;Remark5.9,p190;Table3.5,p151;Theorem1.20 proof,p77; AimPaper Appendix,paragraph immediately beforeTable5; BR21 Proposition1.12;Table1.3,p65,restrictionw2↦v2^8;Table5.4,p196,roww2²; AimPaper Appendix,paragraph immediately beforeTable5: d3(v2^16)=β^5g; Repository AimPaper/main.tex and2412.10879.pdf, frozen baseline. Used only as secondary notation/recall; original BR21 and BMQ checked independently.
% Blueprint note: Detailed source attribution: (Adapted from \cite{BR21}, Theorem 5.18 and Table 5.4, using $\beta^4=g^3$ from Table 3.5; detector transport is a project construction.)
\end{theorem}

\section{Synthetic rigidity and finite quotients}

\begin{definition}[Weight coordinates on the synthetic Adams pages]
\label{sem:Literature:nu-page-coordinates}
\level{coarse}
\uses{sem:AdamsSynthetic:families,frontier:AdamsSynthetic:a5b004b4748d,sem:AdamsSynthetic:shift-coherence,frontier:AdamsSynthetic:f3257e955b3d}
(Adapted from \cite{BHS}, Thm. A.8.)\par{}

For $X$ in the chosen classical object family, weight shift induces isomorphisms
\[E_r^{s,t,w}(\nu X)\cong E_r^{s,t,w-1}(\Sigma^{0,-1}\nu X).\]
For $k\ge0$, the $E_2$ labels in tridegree $(s,t,t-k)$ are the classical labels in bidegree $(s,t)$. Under these coordinates, the map induced by $\lambda:\Sigma^{0,-1}\nu X\to\nu X$ lowers weight by one.
\lean{KIP126.Literature.Standard.nuPage_weightShift,KIP126.Literature.Standard.nuPage_e2,KIP126.Literature.Standard.nuPageData,KIP126.Literature.Standard.pageLambda,KIP126.Literature.Standard.classicalDifferential,KIP126.Literature.Standard.syntheticDifferential}
% Blueprint note: Source verification (2026-10-18): Source adaptation checked. The public shift isomorphism includes integer pages below 2; those pages are frozen to the first page by a project convention, not covered as extra pages in the cited paper. Write the BHS coordinates as (s,k,omega) and the project coordinates as (s,t,w). The translation is k=t-s and omega=w-(t-s); tau is renamed lambda. Hence s-omega=t-w, and the BHS differential degree (r,-1,1) becomes (r,r-1,0). This verifies the mathematical source and its stated regrading, not the outstanding Lean comparison with the fixed project model. Existing Lean proof and transitive-dependency status must remain unchanged. Spectral-sequence page and E-infinity subquotient statements here do not assert convergence to untruncated homotopy groups; BHS Convention A.5 allows an arbitrary spectrum for Theorem A.8 and Corollaries A.9/A.11. Primary references: [Pst](https://arxiv.org/pdf/1803.01804v3): Definition 4.6; Definition 4.27; pages arXiv PDF pp. 39, 45; published pp. 609, 618; arXiv:1803.01804v3, 10 November 2022; numbering also checked in Invent. Math. 232 (2023) [BHS](https://arxiv.org/pdf/1910.14116v3): Theorem A.8; Notation A.6; pages arXiv PDF pp. 71-72; published pp. 306-307; arXiv:1910.14116v3, 10 May 2022; checked against Acta Math. 231 (2023), 205-344 [LWX126](https://arxiv.org/html/2412.10879v2#S3): Theorem 3.6 and Remark 3.7; pages local PDF p. 16; arXiv:2412.10879v2, 22 February 2025; local AimPaper/main.tex and PDF also checked
% Blueprint note: Detailed source attribution: (Adapted from \cite{Pst}, Definitions 4.6 and 4.27, and \cite{BHS}, Theorem A.8; grading as in \cite{LWX126}, Theorem 3.6 and Remark 3.7.)
\end{definition}

\begin{theorem}[The free-lambda comparison range]
\label{sem:Literature:free-lambda-range}
\level{coarse}
\uses{sem:Literature:nu-page-coordinates}
(Adapted from \cite{BHS}, Thm. A.8 and Cor. A.9.)\par{}

For $r\ge2$ and $r-2\le t-w$, there is an isomorphism
\[E_r^{s,t,w}(\nu X)\cong E_r^{s,t}(X).\]
It preserves $E_2$ labels and commutes with $\lambda$ and the differentials. Its inverse gives classical lifts with the same differential and $\lambda$ compatibilities.
\lean{KIP126.Literature.Standard.free_component,KIP126.Literature.Standard.free_e2,KIP126.Literature.Standard.free_lambda,KIP126.Literature.Standard.free_differential,KIP126.Literature.Standard.freeData,KIP126.Literature.Standard.classicalLift,KIP126.Literature.Standard.classicalLift_differential,KIP126.Literature.Standard.classicalLift_lambda}
% Blueprint note: Source verification (2026-10-18): Source adaptation checked. Put a=t-w. In the polynomial-lambda spectral sequence of BHS Theorem A.8, an incoming classical d_i has a source label with lambda exponent a-i+1. All incoming differentials contributing before page r have 2<=i<=r-1, so they all occur once a>=r-2. Outgoing cycles agree under localization. Thus the localization map is an isomorphism for a>=r-2, including the boundary case; this is a derived range, not the literal range printed in Corollary A.9(4). The inverse, E2 normalization, lambda square and differential square are obtained from this same localization map, rather than from unrelated choices of isomorphism. Write the BHS coordinates as (s,k,omega) and the project coordinates as (s,t,w). The translation is k=t-s and omega=w-(t-s); tau is renamed lambda. Hence s-omega=t-w, and the BHS differential degree (r,-1,1) becomes (r,r-1,0). This verifies the mathematical source and its stated regrading, not the outstanding Lean comparison with the fixed project model. Existing Lean proof and transitive-dependency status must remain unchanged. Spectral-sequence page and E-infinity subquotient statements here do not assert convergence to untruncated homotopy groups; BHS Convention A.5 allows an arbitrary spectrum for Theorem A.8 and Corollaries A.9/A.11. Primary references: [BHS](https://arxiv.org/pdf/1910.14116v3): Theorem A.8; Notation A.6; pages arXiv PDF pp. 71-72; published pp. 306-307; arXiv:1910.14116v3, 10 May 2022; checked against Acta Math. 231 (2023), 205-344 [BHS](https://arxiv.org/pdf/1910.14116v3): Corollary A.9(1)-(4), including its final surjectivity sentence; pages arXiv PDF pp. 72-73; published pp. 307-308; arXiv:1910.14116v3, 10 May 2022; checked against Acta Math. 231 (2023), 205-344 [BHS](https://arxiv.org/pdf/1910.14116v3): paragraph immediately before Notation A.7; proof of Theorem A.8; pages arXiv PDF p. 72; published p. 307; arXiv:1910.14116v3, 10 May 2022; checked against Acta Math. 231 (2023), 205-344 [LWX126](https://arxiv.org/html/2412.10879v2#S3): Theorem 3.6 and Remark 3.7; pages local PDF p. 16; arXiv:2412.10879v2, 22 February 2025; local AimPaper/main.tex and PDF also checked
% Blueprint note: Detailed source attribution: (Adapted from \cite{BHS}, Theorem A.8 and Corollary A.9; grading as in \cite{LWX126}, Theorem 3.6 and Remark 3.7.)
\end{theorem}

\begin{definition}[The comparison obtained by inverting lambda]
\label{sem:Literature:lambda-inversion}
\level{coarse}
\uses{sem:Literature:free-lambda-range,sem:AdamsSynthetic:recovery}
(Adapted from \cite{Pst}, Prop. 4.40; \cite{BHS}, Thm. A.8.)\par{}

Inverting $\lambda$ induces a morphism from the synthetic cycle-and-boundary data of $\nu X$ to the classical data. It preserves cycles, boundaries, and $E_2$ labels, commutes with $d_r$ for $r\ge2$, and carries multiplication by $\lambda$ to the identity. For $r-2\le t-w$, its page map is the preceding free-range isomorphism.
\lean{KIP126.Literature.Standard.inversion_data,KIP126.Literature.Standard.inversion_e2,KIP126.Literature.Standard.inversion_differential,KIP126.Literature.Standard.inversion_lambda,KIP126.Literature.Standard.inversionData,KIP126.Literature.Standard.free_inversion}
% Blueprint note: Source verification (2026-10-18): Source adaptation checked. The chosen SSDataMorphism includes all cycle and boundary levels, including the project infinity construction; this is packaging beyond the literal printed spectral-sequence morphism. Write the BHS coordinates as (s,k,omega) and the project coordinates as (s,t,w). The translation is k=t-s and omega=w-(t-s); tau is renamed lambda. Hence s-omega=t-w, and the BHS differential degree (r,-1,1) becomes (r,r-1,0). This verifies the mathematical source and its stated regrading, not the outstanding Lean comparison with the fixed project model. Existing Lean proof and transitive-dependency status must remain unchanged. Spectral-sequence page and E-infinity subquotient statements here do not assert convergence to untruncated homotopy groups; BHS Convention A.5 allows an arbitrary spectrum for Theorem A.8 and Corollaries A.9/A.11. Primary references: [Pst](https://arxiv.org/pdf/1803.01804v3): Proposition 4.36; Theorem 4.37; Proposition 4.40; pages arXiv PDF pp. 47-49; published pp. 621-624; arXiv:1803.01804v3, 10 November 2022; numbering also checked in Invent. Math. 232 (2023) [BHS](https://arxiv.org/pdf/1910.14116v3): paragraph immediately before Notation A.7; proof of Theorem A.8; pages arXiv PDF p. 72; published p. 307; arXiv:1910.14116v3, 10 May 2022; checked against Acta Math. 231 (2023), 205-344 [BHS](https://arxiv.org/pdf/1910.14116v3): Theorem A.8; Notation A.6; pages arXiv PDF pp. 71-72; published pp. 306-307; arXiv:1910.14116v3, 10 May 2022; checked against Acta Math. 231 (2023), 205-344
% Blueprint note: Detailed source attribution: (Adapted from \cite{Pst}, Proposition 4.40, and \cite{BHS}, Notation A.6, the comparison before Notation A.7, and Theorem A.8.)
\end{definition}

\begin{proposition}[The finite quotient E-two strip]
\label{sem:Literature:finite-quotient-labels}
\level{coarse}
\uses{sem:Literature:nu-page-coordinates,sem:AdamsSynthetic:lambda-quotients}
(Adapted from \cite{BHS}, Cor. A.11.)\par{}

For $q>0$,
\[E_2^{s,t,w}(\nu X/\lambda^q)\cong
\begin{cases}E_2^{s,t}(X),&0\le t-w<q,\\0,&\text{otherwise}.\end{cases}\]
In the displayed strip, the quotient projection preserves labels and is an isomorphism on $E_2$.
\lean{KIP126.Literature.Standard.finite_quotient_component,KIP126.Literature.Standard.finite_quotient_component_retained,KIP126.Literature.Standard.finite_quotient_retained,KIP126.Literature.Standard.finite_quotient}
% Blueprint note: Source verification (2026-10-18): Source adaptation checked. The canonical nu-E Adams resolution can be tensored with C-tau^q. Its polynomial E1 terms truncate to exponents 0,...,q-1; the d1 calculation then gives the stated E2 strip. At weight exponent a=t-w, passing to page r inspects outgoing d_i for i<=r-1 whose target exponent is a+i-1. The condition a+r-2<q retains all these targets and all incoming sources, so the actual quotient projection is an isomorphism on that page. This is a deduction from the tower used in BHS Theorem A.8 and Corollary A.11, not a separately numbered safe-range theorem in that paper. Write the BHS coordinates as (s,k,omega) and the project coordinates as (s,t,w). The translation is k=t-s and omega=w-(t-s); tau is renamed lambda. Hence s-omega=t-w, and the BHS differential degree (r,-1,1) becomes (r,r-1,0). This verifies the mathematical source and its stated regrading, not the outstanding Lean comparison with the fixed project model. Existing Lean proof and transitive-dependency status must remain unchanged. Spectral-sequence page and E-infinity subquotient statements here do not assert convergence to untruncated homotopy groups; BHS Convention A.5 allows an arbitrary spectrum for Theorem A.8 and Corollaries A.9/A.11. Primary references: [BHS](https://arxiv.org/pdf/1910.14116v3): Theorem A.8; Notation A.6; pages arXiv PDF pp. 71-72; published pp. 306-307; arXiv:1910.14116v3, 10 May 2022; checked against Acta Math. 231 (2023), 205-344 [BHS](https://arxiv.org/pdf/1910.14116v3): Corollary A.11(1)-(3) and its proof; Notation A.10; pages arXiv PDF p. 73; published p. 308; arXiv:1910.14116v3, 10 May 2022; checked against Acta Math. 231 (2023), 205-344 [LWX126](https://arxiv.org/html/2412.10879v2#S3): Theorem 3.6 and Remark 3.7; pages local PDF p. 16; arXiv:2412.10879v2, 22 February 2025; local AimPaper/main.tex and PDF also checked
% Blueprint note: Detailed source attribution: (Adapted from \cite{BHS}, Theorem A.8, Notation A.10 and Corollary A.11; grading as in \cite{LWX126}, Remark 3.7.)
\end{proposition}

\begin{theorem}[The safe range for finite-quotient comparison]
\label{sem:Literature:finite-quotient-safe}
\level{coarse}
\uses{sem:Literature:finite-quotient-labels,sem:Literature:free-lambda-range}
(Adapted from \cite{BHS}, Thm. A.8 and Cor. A.11.)\par{}

For $q>0$, $r\ge2$, $0\le t-w$, and $t-w+r-2<q$, the quotient projection induces an isomorphism
\[E_r^{s,t,w}(\nu X)\longrightarrow E_r^{s,t,w}(\nu X/\lambda^q).\]
\lean{KIP126.Literature.Standard.finite_quotient_safe}
% Blueprint note: Source verification (2026-10-18): Source adaptation checked. The canonical nu-E Adams resolution can be tensored with C-tau^q. Its polynomial E1 terms truncate to exponents 0,...,q-1; the d1 calculation then gives the stated E2 strip. At weight exponent a=t-w, passing to page r inspects outgoing d_i for i<=r-1 whose target exponent is a+i-1. The condition a+r-2<q retains all these targets and all incoming sources, so the actual quotient projection is an isomorphism on that page. This is a deduction from the tower used in BHS Theorem A.8 and Corollary A.11, not a separately numbered safe-range theorem in that paper. Write the BHS coordinates as (s,k,omega) and the project coordinates as (s,t,w). The translation is k=t-s and omega=w-(t-s); tau is renamed lambda. Hence s-omega=t-w, and the BHS differential degree (r,-1,1) becomes (r,r-1,0). This verifies the mathematical source and its stated regrading, not the outstanding Lean comparison with the fixed project model. Existing Lean proof and transitive-dependency status must remain unchanged. Spectral-sequence page and E-infinity subquotient statements here do not assert convergence to untruncated homotopy groups; BHS Convention A.5 allows an arbitrary spectrum for Theorem A.8 and Corollaries A.9/A.11. Primary references: [BHS](https://arxiv.org/pdf/1910.14116v3): Theorem A.8; Notation A.6; pages arXiv PDF pp. 71-72; published pp. 306-307; arXiv:1910.14116v3, 10 May 2022; checked against Acta Math. 231 (2023), 205-344 [BHS](https://arxiv.org/pdf/1910.14116v3): Corollary A.11(1)-(3) and its proof; Notation A.10; pages arXiv PDF p. 73; published p. 308; arXiv:1910.14116v3, 10 May 2022; checked against Acta Math. 231 (2023), 205-344
% Blueprint note: Detailed source attribution: (Adapted from the finite-quotient Adams tower in \cite{BHS}, Theorem A.8, Notation A.10 and Corollary A.11.)
\end{theorem}

\begin{theorem}[Vanishing windows and inherited degeneration]
\label{sem:Literature:vanishing-and-degeneration}
\level{coarse}
\uses{sem:Literature:nu-page-coordinates,sem:Literature:finite-quotient-labels}
(Adapted from \cite{BHS}, Thm. A.8 and Cor. A.11.)\par{}

For $r\ge2$, one has $E_r^{s,t,w}(\nu X)=0$ when $t<w$. For $q>0$, one has $E_r^{s,t,w}(\nu X/\lambda^q)=0$ when $t-w<0$ or $q\le t-w$. If the classical Adams spectral sequence of $X$ degenerates at $N$, then the synthetic Adams spectral sequence of $\nu X$ degenerates at $N$.
\lean{KIP126.Literature.Standard.above_diagonal,KIP126.Literature.Standard.finite_quotient_vanishing,KIP126.Literature.Standard.degeneration}
% Blueprint note: Source verification (2026-10-18): Partial source correspondence. No strong-convergence conclusion is asserted by degeneration alone. The displayed mathematical assertion is sourced for N>=2. The bound Lean degeneration interface quantifies every integer N; since DegeneratesAt asks for every differential r>=N, the N<2 portion additionally requires a below-first-page normalization for the actual synthetic family. BHS Theorem A.8 does not provide that project-specific portion. The canonical classical construction does set r<2 differentials to zero in Internal/Source/AdamsSequence.lean:207, but no matching synthetic-family normalization was found in this research. This is an unverified scope endpoint, not a verified counterexample or authorization to edit an interface. Write the BHS coordinates as (s,k,omega) and the project coordinates as (s,t,w). The translation is k=t-s and omega=w-(t-s); tau is renamed lambda. Hence s-omega=t-w, and the BHS differential degree (r,-1,1) becomes (r,r-1,0). This verifies the mathematical source and its stated regrading, not the outstanding Lean comparison with the fixed project model. Existing Lean proof and transitive-dependency status must remain unchanged. Spectral-sequence page and E-infinity subquotient statements here do not assert convergence to untruncated homotopy groups; BHS Convention A.5 allows an arbitrary spectrum for Theorem A.8 and Corollaries A.9/A.11. Primary references: [BHS](https://arxiv.org/pdf/1910.14116v3): Theorem A.8; Notation A.6; pages arXiv PDF pp. 71-72; published pp. 306-307; arXiv:1910.14116v3, 10 May 2022; checked against Acta Math. 231 (2023), 205-344 [BHS](https://arxiv.org/pdf/1910.14116v3): Corollary A.9(1)-(4), including its final surjectivity sentence; pages arXiv PDF pp. 72-73; published pp. 307-308; arXiv:1910.14116v3, 10 May 2022; checked against Acta Math. 231 (2023), 205-344 [BHS](https://arxiv.org/pdf/1910.14116v3): Corollary A.11(1)-(3) and its proof; Notation A.10; pages arXiv PDF p. 73; published p. 308; arXiv:1910.14116v3, 10 May 2022; checked against Acta Math. 231 (2023), 205-344
% Blueprint note: Detailed source attribution: (Adapted from \cite{BHS}, Theorem A.8 and Corollaries A.9(1) and A.11(1),(3); the extension to $N<2$ remains unverified.)
\end{theorem}

\section{Infinite pages and weight maps}

\begin{definition}[Classical cycles, boundaries, and common permanent representatives]
\label{sem:Literature:cycle-quotients}
\level{coarse}
\uses{sem:AdamsSynthetic:families,sem:SpectralSequence:cycles-boundaries,sem:SpectralSequence:permanent}
(Project construction; cf.\ \cite{LWX126}, Notation 3.10.)\par{}

Let $Z_r^{s,t}(X)$ and $B_r^{s,t}(X)$ be the classical cycle and boundary subgroups of $E_2^{s,t}(X)$, and let $Z_\infty=\bigcap_r Z_r$. A classical permanent representative is an element of $Z_\infty$; zero is allowed. The cycle and permanent quotients are the corresponding subquotients of $E_2^{s,t}(X)$.
\lean{KIP126.Literature.Standard.cycles,KIP126.Literature.Standard.boundaries,KIP126.Literature.Standard.permanentCycles,KIP126.Literature.Standard.CycleQuotient,KIP126.Literature.Standard.PermanentQuotient,KIP126.Literature.Standard.quotientMap,KIP126.Literature.Standard.ClassicalPermanentCycle}
% Blueprint note: Source verification (2026-10-18): Project choice or construction. Do not describe ClassicalPermanentCycle as a nonzero permanent cycle or as a convergent detected class. Write the BHS coordinates as (s,k,omega) and the project coordinates as (s,t,w). The translation is k=t-s and omega=w-(t-s); tau is renamed lambda. Hence s-omega=t-w, and the BHS differential degree (r,-1,1) becomes (r,r-1,0). This verifies the mathematical source and its stated regrading, not the outstanding Lean comparison with the fixed project model. Existing Lean proof and transitive-dependency status must remain unchanged. Spectral-sequence page and E-infinity subquotient statements here do not assert convergence to untruncated homotopy groups; BHS Convention A.5 allows an arbitrary spectrum for Theorem A.8 and Corollaries A.9/A.11. Primary references: [BHS](https://arxiv.org/pdf/1910.14116v3): Notation A.7; pages arXiv PDF p. 72; published p. 307; arXiv:1910.14116v3, 10 May 2022; checked against Acta Math. 231 (2023), 205-344 [LWX126](https://arxiv.org/html/2412.10879v2#S3): Notation 3.10; pages local PDF pp. 16-17; arXiv:2412.10879v2, 22 February 2025; local AimPaper/main.tex and PDF also checked
% Blueprint note: Detailed source attribution: (Adapted from \cite{BHS}, Notation A.7, and \cite{LWX126}, Notation 3.10; project subquotient definitions.)
\end{definition}

\begin{proposition}[The infinite page of nu X and its labels]
\label{sem:Literature:infinite-nu-window}
\level{coarse}
\uses{sem:Literature:cycle-quotients,sem:Literature:nu-page-coordinates,sem:Literature:vanishing-and-degeneration,paper:3.6}
(Adapted from \cite{BHS}, Cor. A.9.)\par{}

For $X$ in the chosen classical object family,
\[E_\infty^{s,t,w}(\nu X)\cong
\begin{cases}Z_\infty^{s,t}(X)/B_{1+t-w}^{s,t}(X),&w\le t,\\0,&t<w.\end{cases}\]
This is a linear isomorphism carrying the infinite representative of an $E_2$ label to its permanent-cycle quotient class.
\lean{KIP126.Literature.Standard.eInfinity_nu,KIP126.Literature.Standard.eInfinity_nu_labels,KIP126.Literature.Standard.infinityNuWindow,KIP126.Literature.Standard.infinityNuLinear}
% Blueprint note: Source verification (2026-10-18): Source adaptation checked. BHS Corollary A.9 gives the untruncated E-infinity subquotient, and Corollary A.11 gives the finite quotient. All formula indices transform by s-omega=t-w. The finite formula holds for q>=1 in BHS, while AimPaper Proposition 3.12 explicitly states q>=2. Additive equivalences are viewed as Z-linear without a change of scalar structure. Matching a specific E2 representative to its quotient class requires the canonical label-compatible presentation, not an arbitrary abstract isomorphism. Zero labels are permitted in subgroup and quotient statements. Write the BHS coordinates as (s,k,omega) and the project coordinates as (s,t,w). The translation is k=t-s and omega=w-(t-s); tau is renamed lambda. Hence s-omega=t-w, and the BHS differential degree (r,-1,1) becomes (r,r-1,0). This verifies the mathematical source and its stated regrading, not the outstanding Lean comparison with the fixed project model. Existing Lean proof and transitive-dependency status must remain unchanged. Spectral-sequence page and E-infinity subquotient statements here do not assert convergence to untruncated homotopy groups; BHS Convention A.5 allows an arbitrary spectrum for Theorem A.8 and Corollaries A.9/A.11. Primary references: [BHS](https://arxiv.org/pdf/1910.14116v3): Corollary A.9(1)-(4), including its final surjectivity sentence; pages arXiv PDF pp. 72-73; published pp. 307-308; arXiv:1910.14116v3, 10 May 2022; checked against Acta Math. 231 (2023), 205-344 [LWX126](https://arxiv.org/html/2412.10879v2#S3): Proposition 3.11; pages local PDF p. 17; arXiv:2412.10879v2, 22 February 2025; local AimPaper/main.tex and PDF also checked
% Blueprint note: Detailed source attribution: (Adapted from \cite{BHS}, Corollary A.9, and \cite{LWX126}, Proposition 3.11.)
\end{proposition}

\begin{proposition}[The infinite page of a finite lambda quotient]
\label{sem:Literature:infinite-finite-window}
\level{coarse}
\uses{sem:Literature:cycle-quotients,sem:Literature:finite-quotient-labels,sem:Literature:vanishing-and-degeneration,paper:3.6}
(Adapted from \cite{BHS}, Cor. A.11.)\par{}

For $q>0$,
\[E_\infty^{s,t,w}(\nu X/\lambda^q)\cong
\begin{cases}Z_{q-t+w}^{s,t}(X)/B_{1+t-w}^{s,t}(X),&0\le t-w<q,\\0,&\text{otherwise}.\end{cases}\]
This is a linear isomorphism compatible with labels and the quotient projection.
\lean{KIP126.Literature.Standard.eInfinity_finite,KIP126.Literature.Standard.eInfinity_finite_labels,KIP126.Literature.Standard.infinityFiniteWindow,KIP126.Literature.Standard.infinityFiniteLinear}
% Blueprint note: Source verification (2026-10-18): Source adaptation checked. BHS Corollary A.9 gives the untruncated E-infinity subquotient, and Corollary A.11 gives the finite quotient. All formula indices transform by s-omega=t-w. The finite formula holds for q>=1 in BHS, while AimPaper Proposition 3.12 explicitly states q>=2. Additive equivalences are viewed as Z-linear without a change of scalar structure. Matching a specific E2 representative to its quotient class requires the canonical label-compatible presentation, not an arbitrary abstract isomorphism. Zero labels are permitted in subgroup and quotient statements. Write the BHS coordinates as (s,k,omega) and the project coordinates as (s,t,w). The translation is k=t-s and omega=w-(t-s); tau is renamed lambda. Hence s-omega=t-w, and the BHS differential degree (r,-1,1) becomes (r,r-1,0). This verifies the mathematical source and its stated regrading, not the outstanding Lean comparison with the fixed project model. Existing Lean proof and transitive-dependency status must remain unchanged. Spectral-sequence page and E-infinity subquotient statements here do not assert convergence to untruncated homotopy groups; BHS Convention A.5 allows an arbitrary spectrum for Theorem A.8 and Corollaries A.9/A.11. Primary references: [BHS](https://arxiv.org/pdf/1910.14116v3): Corollary A.11(1)-(3) and its proof; Notation A.10; pages arXiv PDF p. 73; published p. 308; arXiv:1910.14116v3, 10 May 2022; checked against Acta Math. 231 (2023), 205-344 [LWX126](https://arxiv.org/html/2412.10879v2#S3): Proposition 3.12; pages local PDF p. 17; arXiv:2412.10879v2, 22 February 2025; local AimPaper/main.tex and PDF also checked [LWX126](https://arxiv.org/html/2412.10879v2#S3): Remark 3.14; pages local PDF p. 17; arXiv:2412.10879v2, 22 February 2025; local AimPaper/main.tex and PDF also checked
% Blueprint note: Detailed source attribution: (Adapted from \cite{BHS}, Corollary A.11, and \cite{LWX126}, Proposition 3.12 and Remark 3.14.)
\end{proposition}

\begin{proposition}[Shift comparison and naturality at infinity]
\label{sem:Literature:infinite-shift}
\level{coarse}
\uses{sem:AdamsSynthetic:ess-shifts,sem:Literature:infinite-nu-window,sem:Literature:infinite-finite-window}
(Project construction; cf.\ \cite{LWX126}, Rem. 3.14.)\par{}

For $k\ge0$, the shift $(0,-k)$ induces a filtered comparison and an $E_\infty$ comparison, natural in morphisms between objects in the chosen synthetic family. Composing with the map $\lambda^k$ gives the weight-lowering maps on the untruncated and finite-quotient infinite pages. At $k=0$ these maps are identities.
\lean{KIP126.Literature.Standard.infinityShift_comparison,KIP126.Literature.Standard.infinityShift_natural,KIP126.Literature.Standard.infinityShiftData,KIP126.Literature.Standard.infinityLambda,KIP126.Literature.Standard.infinityFiniteLambda,KIP126.Literature.Standard.infinityLambda_zero,KIP126.Literature.Standard.infinityFiniteLambda_zero}
% Blueprint note: Source verification (2026-10-18): Project choice or construction. This is not an independently quoted literature theorem. The project proves naturality from the chosen filtered comparison, so source attribution should identify the construction. Write the BHS coordinates as (s,k,omega) and the project coordinates as (s,t,w). The translation is k=t-s and omega=w-(t-s); tau is renamed lambda. Hence s-omega=t-w, and the BHS differential degree (r,-1,1) becomes (r,r-1,0). This verifies the mathematical source and its stated regrading, not the outstanding Lean comparison with the fixed project model. Existing Lean proof and transitive-dependency status must remain unchanged. Spectral-sequence page and E-infinity subquotient statements here do not assert convergence to untruncated homotopy groups; BHS Convention A.5 allows an arbitrary spectrum for Theorem A.8 and Corollaries A.9/A.11. Primary references: [Pst](https://arxiv.org/pdf/1803.01804v3): Definition 4.6; Definition 4.27; pages arXiv PDF pp. 39, 45; published pp. 609, 618; arXiv:1803.01804v3, 10 November 2022; numbering also checked in Invent. Math. 232 (2023) [BHS](https://arxiv.org/pdf/1910.14116v3): Theorem A.8; Notation A.6; pages arXiv PDF pp. 71-72; published pp. 306-307; arXiv:1910.14116v3, 10 May 2022; checked against Acta Math. 231 (2023), 205-344 [LWX126](https://arxiv.org/html/2412.10879v2#S3): Remark 3.14; pages local PDF p. 17; arXiv:2412.10879v2, 22 February 2025; local AimPaper/main.tex and PDF also checked
% Blueprint note: Detailed source attribution: (Adapted from \cite{Pst}, Definitions 4.6 and 4.27, \cite{BHS}, Notation A.6, and \cite{LWX126}, Remark 3.14; project shift comparisons.)
\end{proposition}

\begin{proposition}[Lambda on infinite-page quotient coordinates]
\label{sem:Literature:infinite-lambda}
\level{coarse}
\uses{sem:Literature:infinite-shift,sem:Literature:infinite-nu-window,sem:Literature:infinite-finite-window}
(Adapted from \cite{LWX126}, Rem. 3.14.)\par{}

Under the infinite-page presentations, multiplication by $\lambda^k$ induces the natural map between classical cycle/boundary quotients. For $\nu X$, assume $w\le t$. For finite quotients, assume $k<q$ and $0\le t-w<q-k$; the source has length $q-k$ and weight $w$, and the target has length $q$ and weight $w-k$.
\lean{KIP126.Literature.Standard.eInfinity_lambda_nu,KIP126.Literature.Standard.eInfinity_lambda_finite}
% Blueprint note: Source verification (2026-10-18): Source adaptation checked. BHS Corollary A.9 gives the untruncated E-infinity subquotient, and Corollary A.11 gives the finite quotient. All formula indices transform by s-omega=t-w. The finite formula holds for q>=1 in BHS, while AimPaper Proposition 3.12 explicitly states q>=2. Additive equivalences are viewed as Z-linear without a change of scalar structure. Matching a specific E2 representative to its quotient class requires the canonical label-compatible presentation, not an arbitrary abstract isomorphism. Zero labels are permitted in subgroup and quotient statements. Write the BHS coordinates as (s,k,omega) and the project coordinates as (s,t,w). The translation is k=t-s and omega=w-(t-s); tau is renamed lambda. Hence s-omega=t-w, and the BHS differential degree (r,-1,1) becomes (r,r-1,0). This verifies the mathematical source and its stated regrading, not the outstanding Lean comparison with the fixed project model. Existing Lean proof and transitive-dependency status must remain unchanged. Spectral-sequence page and E-infinity subquotient statements here do not assert convergence to untruncated homotopy groups; BHS Convention A.5 allows an arbitrary spectrum for Theorem A.8 and Corollaries A.9/A.11. Primary references: [BHS](https://arxiv.org/pdf/1910.14116v3): Corollary A.9(1)-(4), including its final surjectivity sentence; pages arXiv PDF pp. 72-73; published pp. 307-308; arXiv:1910.14116v3, 10 May 2022; checked against Acta Math. 231 (2023), 205-344 [BHS](https://arxiv.org/pdf/1910.14116v3): Corollary A.11(1)-(3) and its proof; Notation A.10; pages arXiv PDF p. 73; published p. 308; arXiv:1910.14116v3, 10 May 2022; checked against Acta Math. 231 (2023), 205-344 [LWX126](https://arxiv.org/html/2412.10879v2#S3): Remark 3.14; pages local PDF p. 17; arXiv:2412.10879v2, 22 February 2025; local AimPaper/main.tex and PDF also checked
% Blueprint note: Detailed source attribution: (Adapted from \cite{BHS}, Corollaries A.9 and A.11, and \cite{LWX126}, Remark 3.14.)
\end{proposition}

\begin{proposition}[Restriction between infinite-page windows]
\label{sem:Literature:infinite-restriction}
\level{coarse}
\uses{sem:Literature:infinite-nu-window,sem:Literature:infinite-finite-window}
(Adapted from \cite{LWX126}, Rem. 3.14.)\par{}

For $0<i\le j$ and $0\le t-w<i$, restriction $\nu X/\lambda^j\to\nu X/\lambda^i$ induces the natural map between the cycle/boundary quotients. Projection from $\nu X$ to a positive-length finite quotient carries a permanent-cycle quotient class to its finite-cycle quotient class in the common strip.
\lean{KIP126.Literature.Standard.eInfinity_rho_finite,KIP126.Literature.Standard.eInfinity_rho_nu}
% Blueprint note: Source verification (2026-10-18): Source adaptation checked. BHS Corollary A.9 gives the untruncated E-infinity subquotient, and Corollary A.11 gives the finite quotient. All formula indices transform by s-omega=t-w. The finite formula holds for q>=1 in BHS, while AimPaper Proposition 3.12 explicitly states q>=2. Additive equivalences are viewed as Z-linear without a change of scalar structure. Matching a specific E2 representative to its quotient class requires the canonical label-compatible presentation, not an arbitrary abstract isomorphism. Zero labels are permitted in subgroup and quotient statements. Write the BHS coordinates as (s,k,omega) and the project coordinates as (s,t,w). The translation is k=t-s and omega=w-(t-s); tau is renamed lambda. Hence s-omega=t-w, and the BHS differential degree (r,-1,1) becomes (r,r-1,0). This verifies the mathematical source and its stated regrading, not the outstanding Lean comparison with the fixed project model. Existing Lean proof and transitive-dependency status must remain unchanged. Spectral-sequence page and E-infinity subquotient statements here do not assert convergence to untruncated homotopy groups; BHS Convention A.5 allows an arbitrary spectrum for Theorem A.8 and Corollaries A.9/A.11. Primary references: [BHS](https://arxiv.org/pdf/1910.14116v3): Corollary A.11(1)-(3) and its proof; Notation A.10; pages arXiv PDF p. 73; published p. 308; arXiv:1910.14116v3, 10 May 2022; checked against Acta Math. 231 (2023), 205-344 [BHS](https://arxiv.org/pdf/1910.14116v3): Corollary A.9(1)-(4), including its final surjectivity sentence; pages arXiv PDF pp. 72-73; published pp. 307-308; arXiv:1910.14116v3, 10 May 2022; checked against Acta Math. 231 (2023), 205-344 [LWX126](https://arxiv.org/html/2412.10879v2#S3): Remark 3.14; pages local PDF p. 17; arXiv:2412.10879v2, 22 February 2025; local AimPaper/main.tex and PDF also checked
% Blueprint note: Detailed source attribution: (Adapted from \cite{BHS}, Corollary A.11 and its proof, and \cite{LWX126}, Remark 3.14.)
\end{proposition}

\section{Cofibers, filtered lifts, and normalized triangles}

\begin{theorem}[Synthetic cofibers and full lifts of maps]
\label{sem:Literature:cofiber-and-lift}
\level{coarse}
\uses{sem:AdamsSynthetic:nu-exactness,sem:AdamsSynthetic:filtration-comparison}
(Adapted from \cite{Pst}, Lem. 4.23; \cite{BHS}, Lem. 9.15.)\par{}

For a classical cofiber triangle, its image under $\nu$ is a synthetic cofiber triangle if and only if its mod-two homology sequence is short exact. If $f$ has Adams tower filtration at least $k$, it has a synthetic lift whose composite with $\lambda^k$ is $\nu f$. A positive-filtration connecting map admits such a lift as the connecting arrow of the lifted triangle. Two full lifts may differ by $\lambda$ torsion.
\lean{KIP126.Literature.Standard.nu_cofiber,KIP126.Literature.Standard.lift,KIP126.Literature.Standard.triangle_lift}
% Blueprint note: Source verification (2026-10-18): Source adaptation checked. BHS factors a map of filtration k. The interface uses the actual target Adams-tower lower-bound predicate and negative-target shifts; identifying that predicate with source Adams filtration is part of model recognition. The interface existentially chooses a connecting arrow and its distinguished triangle under positive filtration and homology exactness; its full-lift comparison is only modulo lambda torsion, as in AimPaper Remark 3.18. Specialization to mod-two homology and the project’s precise nu and coefficient functors; no finiteness, connectivity, or completeness premise is introduced. Primary references: [BHS](https://arxiv.org/pdf/1910.14116v3): Lemma 9.15 and its proof; pages 34–35; arXiv:1910.14116v3, 10 May 2022 [LWX126](https://arxiv.org/abs/2412.10879): Propositions 3.16–3.17; Remark 3.18; pages 18–19; Repository-frozen AimPaper/main.tex and AimPaper/2412.10879.pdf [Pst](https://arxiv.org/pdf/1803.01804v3): Lemma 4.23; pages 43; arXiv:1803.01804v3, 10 November 2022
% Blueprint note: Detailed source attribution: (Adapted from \cite{Pst}, Lemma 4.23, and \cite{BHS}, Lemma 9.15; see \cite{LWX126}, Propositions 3.16--3.17 and Remark 3.18.)
\end{theorem}

\begin{definition}[Normalized arrows and their actual cofiber comparison]
\label{sem:Literature:normalized-cofiber}
\level{coarse}
\uses{sem:Literature:cofiber-and-lift,sem:AdamsSynthetic:normalized-map,frontier:AdamsSynthetic:d505f4e20090}
(Adapted from \cite{LWX126}, Prop. 3.20.)\par{}

For a triangle in the chosen classical object family, let its normalized arrows be obtained from synthetic lifts by shifting through their filtration exponents. The cofiber triangle of the first normalized arrow is distinguished and has a cofiber comparison. Call the normalized arrows compatible when the triangle formed by all three arrows is distinguished.
% Blueprint note: An exponent sum of one is not bound to a general theorem proving compatibility of all three prescribed normalized arrows. The unshifted existence statement only gives some distinguished triangle on the three objects.
\lean{KIP126.Literature.Standard.normalizedTriangle,KIP126.Literature.Standard.NormalizedTriangleCompatible,KIP126.Literature.Standard.normalized_cofiber,KIP126.Literature.Standard.normalizedCofiberTriangle,KIP126.Literature.Standard.normalizedCofiberTriangle_distinguished,KIP126.Literature.Standard.unshifted}
% Blueprint note: Source verification (2026-10-18): Partial source correspondence. Only the conditional/coordinated source construction is verified. The all-triangles selected cofiber comparison and unshifted triangle are not supplied by Proposition 3.20; do not describe either as an exact restatement. Primary references: [LWX126](https://arxiv.org/abs/2412.10879): Notation 3.19; Proposition 3.20; pages 19; Repository-frozen AimPaper/main.tex and AimPaper/2412.10879.pdf
% Blueprint note: Detailed source attribution: (Adapted from \cite{LWX126}, Notation 3.19 and Proposition 3.20; comparison with the selected normalized maps remains to be verified.)
\end{definition}

\begin{proposition}[The normalized cofiber triangle of nu]
\label{sem:Literature:nu-cofiber}
\level{coarse}
\uses{sem:Literature:hopf-classes,sem:Literature:normalized-cofiber}
(Adapted from \cite{LWX126}, Ex. 6.19.)\par{}

For the classical triangle $S^3\xrightarrow{\nu}S^0\to C\nu\to S^4$, the normalized exponent of $\nu$ is one and those of the bottom- and top-cell maps are zero. The normalized $\nu$ has first-quotient label $h_2$, and the triangle of normalized arrows is distinguished.
% Blueprint note: The compatibility of this particular normalized triangle retains its separate comparison obligation.
\lean{KIP126.Literature.Standard.nuCofiber_nu_exponent,KIP126.Literature.Standard.nuCofiber_bottom_exponent,KIP126.Literature.Standard.nuCofiber_top_exponent,KIP126.Literature.Standard.nuCofiber_normalized_label,KIP126.Literature.Standard.nuCofiber_triangle,KIP126.Literature.Standard.nuCofiber_exponent_sum,KIP126.Literature.Standard.nuCofiber_distinguished,KIP126.Literature.Standard.normalizedNu}
% Blueprint note: Source verification (2026-10-18): Partial source correspondence. Example 6.19 supplies the exponent values and the Ext h2 operation. The first-quotient label of normalizedNu and distinguishedness of the three already selected maps are separate obligations. Primary references: [LWX126](https://arxiv.org/abs/2412.10879): Example 6.19; Proposition 3.20; pages 38–39; 19; Repository-frozen AimPaper/main.tex and AimPaper/2412.10879.pdf [LWX126](https://arxiv.org/abs/2412.10879): Notation 3.19; Proposition 3.20; pages 19; Repository-frozen AimPaper/main.tex and AimPaper/2412.10879.pdf
% Blueprint note: Detailed source attribution: (Adapted from \cite{LWX126}, Example 6.19 and Proposition 3.20; the selected-map comparisons remain explicit obligations.)
\end{proposition}

\section{Bockstein lifts and completion hypotheses}

\begin{theorem}[Finite lifts and the Bockstein differential]
\label{sem:Literature:finite-bockstein}
\level{coarse}
\uses{sem:AdamsSynthetic:lambda-quotients,sem:AdamsSynthetic:e2-comparison,sem:Literature:finite-quotient-labels}
(Adapted from \cite{BHS}, Thm. A.1(1) and Cor. A.11.)\par{}

Let $q>0$ and $x\in E_2^{s,t}(X)$. The label $x$ reaches $E_{q+1}$ if and only if its first-quotient class lifts to $\nu X/\lambda^q$. Here reaching a page allows zero and boundaries. A suitable lift has Bockstein label representing $d_{q+1}x$ in bidegree $(s+q+1,t+q)$, where the Bockstein is the boundary of this quotient tower.
\lean{KIP126.Literature.Standard.bhs_finite_lift,KIP126.Literature.Standard.bhs_bockstein}
% Blueprint note: Source verification (2026-10-18): Source adaptation checked. Finite-quotient result: write q for BHS r, identify tau with lambda, and use stem k=t-s. Reaching E_(q+1) means d_2 through d_q vanish; zero and earlier boundaries are allowed. Unlike the untruncated conclusion, the finite case extends to arbitrary spectra via Corollary A.11 and Lemma A.14, as in the proof of Theorem A.1(1). Over F_2 the sign in A.1(1c) disappears. The mathematical source has been checked; identification with the project's specified coefficient object, Adams towers, quotient maps and realization remains part of the existing formal comparison obligations. No Lean proof obligation is discharged by this citation. Primary references: [BHS](https://arxiv.org/pdf/1910.14116v3): Theorem A.1(1a)–(1c); Corollary A.11; Lemma A.14; proof of Theorem 9.19(1); pages PDF pp. 70, 73–75; arXiv:1910.14116v3, 10 May 2022; final prepublication version, not published pagination [LWX126](https://arxiv.org/abs/2412.10879): Theorem 3.8(1)–(2); pages PDF p. 16; Repository AimPaper/main.tex and AimPaper/2412.10879.pdf; exact local source snapshot
% Blueprint note: Detailed source attribution: (Adapted from \cite{BHS}, Theorem A.1(1), Corollary A.11 and Lemma A.14; \cite{LWX126}, Theorem 3.8.)
\end{theorem}

\begin{theorem}[Classical and synthetic differentials with common labels]
\label{sem:Literature:differential-rigidity}
\level{coarse}
\uses{sem:Literature:free-lambda-range,sem:Literature:nu-page-coordinates}
(Adapted from \cite{BHS}, Thm. A.8(3).)\par{}

For $r\ge2$, the classical relation $d_r(x)=y$ is equivalent to the synthetic labeled relation $d_r(x)=\lambda^{r-1}y$. More generally,
\[d_r(\lambda^k x)=\lambda^{k+r-1}y.\]
After a shift $(0,a)$, both sides have weight $t+a-k$.
\lean{KIP126.Literature.Standard.bhs_differentials}
% Blueprint note: Source verification (2026-10-18): Source adaptation checked. Tau is renamed lambda, the third grading is changed as specified in Remark 3.7, and the formula is extended by lambda-linearity and the shift (0,a). The displayed source and target both have weight t+a-k. The Lean relation HasDifferential requires representatives of both E_2 labels on E_r; it permits a zero target. The mathematical source has been checked; identification with the project's specified coefficient object, Adams towers, quotient maps and realization remains part of the existing formal comparison obligations. No Lean proof obligation is discharged by this citation. Primary references: [BHS](https://arxiv.org/pdf/1910.14116v3): Theorem A.8(2)–(3); pages PDF p. 72; arXiv:1910.14116v3, 10 May 2022; final prepublication version, not published pagination [LWX126](https://arxiv.org/abs/2412.10879): Theorem 3.6 and Remark 3.7; pages PDF p. 16; Repository AimPaper/main.tex and AimPaper/2412.10879.pdf; exact local source snapshot
% Blueprint note: Detailed source attribution: (Adapted from \cite{BHS}, Theorem A.8(3); \cite{LWX126}, Theorem 3.6 and Remark 3.7.)
\end{theorem}

\begin{definition}[The complete-source hypotheses for untruncated BHS]
\label{sem:Literature:bhs-applicability}
\level{coarse}
(Project construction; cf.\ \cite{BHS}, Def. A.12.)\par{}

An object satisfies the complete-source hypotheses if countable products exist, the Milnor map $1-\mathrm{shift}$ on the product of residual Adams-tower stages is invertible, and the tower admits a completion witness. Require also that the homotopy filtration is separated in every degree and that the canonical convergence identification holds for this tower.
% Blueprint note: The complete-source hypotheses include derived-limit information and are not inferred merely from an $E_\infty$ isomorphism.
\lean{KIP126.Literature.Standard.BHSObjectApplicability}
% Blueprint note: Source verification (2026-10-18): Project choice or construction. This is a project definition. Its Milnor 1-minus-shift is the actual map on the product of residual stages, so invertibility encodes vanishing of the sequential homotopy limit, including derived-limit information. Its completion witnesses, separatedness, and associated-graded convergence are retained separately. BHS does not define this Lean structure or assert that every selected object satisfies it. Primary references: [BHS](https://arxiv.org/pdf/1910.14116v3): Definitions 9.16–9.17; pages PDF p. 35; arXiv:1910.14116v3, 10 May 2022; final prepublication version, not published pagination [BHS](https://arxiv.org/pdf/1910.14116v3): Definition A.12; Propositions A.13 and A.16; pages PDF pp. 73–76; arXiv:1910.14116v3, 10 May 2022; final prepublication version, not published pagination
% Blueprint note: Detailed source attribution: (Adapted from \cite{BHS}, Definitions 9.16--9.17 and A.12, and Propositions A.13 and A.16; project Adams-tower formulation.)
\end{definition}

\begin{theorem}[Permanent lifts and filtration on a complete source]
\label{sem:Literature:complete-permanent-lifts}
\level{coarse}
(Adapted from \cite{BHS}, Thm. A.1 and Cor. A.19.)\par{}

Assume the complete-source hypotheses. A classical label has a common infinite-cycle representative if and only if it is the first-quotient image of an untruncated synthetic class. For a class of bidegree $(m,w)$ and $w-m\le s$, Adams filtration at least $s$ is equivalent to divisibility by $\lambda^{m+s-w}$ from weight $m+s$.
\lean{KIP126.Literature.Standard.bhs_permanent_lift,KIP126.Literature.Standard.bhs_filtration_lambda}
% Blueprint note: Source verification (2026-10-18): Source adaptation checked. BHS assumes E-nilpotent completeness and a strongly convergent E-based Adams spectral sequence (Definitions 9.16–9.17). BHSObjectApplicability is the project's explicit Adams-tower encoding of these hypotheses, not a quoted definition from BHS. The mathematical source has been checked; identification with the project's specified coefficient object, Adams towers, quotient maps and realization remains part of the existing formal comparison obligations. No Lean proof obligation is discharged by this citation. Primary references: [BHS](https://arxiv.org/pdf/1910.14116v3): Theorem A.1, permanent-cycle paragraph preceding (2a); proof of Theorem 9.19(2)–(4); pages PDF pp. 70, 76–77; arXiv:1910.14116v3, 10 May 2022; final prepublication version, not published pagination [BHS](https://arxiv.org/pdf/1910.14116v3): Corollary A.19; pages PDF p. 76; arXiv:1910.14116v3, 10 May 2022; final prepublication version, not published pagination
% Blueprint note: Detailed source attribution: (Adapted from \cite{BHS}, Theorem A.1, the permanent-cycle paragraph, and Corollary A.19.)
\end{theorem}

\begin{theorem}[Lifetime of a lift on a complete source]
\label{sem:Literature:complete-lifetime}
\level{coarse}
(Adapted from \cite{BHS}, Thm. A.1(2a).)\par{}

Assume the complete-source hypotheses. Let $r\ge1$, and let a synthetic class $a$ have first-quotient label $x$. If $x$ has a common permanent representative and survives nontrivially to $E_{r+1}$, then $\lambda^{r-1}a\ne0$.
\lean{KIP126.Literature.Standard.bhs_realization_lifetime}
% Blueprint note: Source verification (2026-10-18): Source adaptation checked. BHS assumes E-nilpotent completeness and a strongly convergent E-based Adams spectral sequence (Definitions 9.16–9.17). BHSObjectApplicability is the project's explicit Adams-tower encoding of these hypotheses, not a quoted definition from BHS. The mathematical source has been checked; identification with the project's specified coefficient object, Adams towers, quotient maps and realization remains part of the existing formal comparison obligations. No Lean proof obligation is discharged by this citation. Primary references: [BHS](https://arxiv.org/pdf/1910.14116v3): Theorem A.1(2a), equivalently Theorem 9.19(2a); pages PDF pp. 70, 77 (also p. 36); arXiv:1910.14116v3, 10 May 2022; final prepublication version, not published pagination
% Blueprint note: Detailed source attribution: (Adapted from \cite{BHS}, Theorem A.1(2a), equivalently Theorem 9.19(2a).)
\end{theorem}

\begin{theorem}[Detection and prescribed realization on a complete source]
\label{sem:Literature:complete-detection}
\level{coarse}
(Adapted from \cite{BHS}, Thm. A.1(2b),(3b).)\par{}

Assume the complete-source hypotheses. The realization of a synthetic lift of a nonzero permanent label $x$ is detected by $x$ in the classical tower. Conversely, if a classical class $a$ is detected by such an $x$, there is a synthetic class with first-quotient label $x$ and realization $a$.
\lean{KIP126.Literature.Standard.bhs_realization_detection,KIP126.Literature.Standard.bhs_realization_prescribed_lift}
% Blueprint note: Source verification (2026-10-18): Source adaptation checked. BHS assumes E-nilpotent completeness and a strongly convergent E-based Adams spectral sequence (Definitions 9.16–9.17). BHSObjectApplicability is the project's explicit Adams-tower encoding of these hypotheses, not a quoted definition from BHS. The mathematical source has been checked; identification with the project's specified coefficient object, Adams towers, quotient maps and realization remains part of the existing formal comparison obligations. No Lean proof obligation is discharged by this citation. Primary references: [BHS](https://arxiv.org/pdf/1910.14116v3): Theorem A.1(2b) and (3b), equivalently Theorem 9.19(2b),(3b); pages PDF pp. 70, 77 (also p. 36); arXiv:1910.14116v3, 10 May 2022; final prepublication version, not published pagination
% Blueprint note: Detailed source attribution: (Adapted from \cite{BHS}, Theorem A.1(2b),(3b), equivalently Theorem 9.19(2b),(3b).)
\end{theorem}

\begin{theorem}[A boundary has a suitably torsion lift]
\label{sem:Literature:complete-boundary-lift}
\level{coarse}
(Adapted from \cite{BHS}, Thm. A.1(3a).)\par{}

Assume the complete-source hypotheses. If $r\ge2$ and a classical label $y$ has a common infinite-cycle representative and is hit on page $r$, then $y$ has a synthetic lift $a$ with $\lambda^{r-1}a=0$.
\lean{KIP126.Literature.Standard.bhs_realization_boundary_lift}
% Blueprint note: Source verification (2026-10-18): Source adaptation checked. BHS assumes E-nilpotent completeness and a strongly convergent E-based Adams spectral sequence (Definitions 9.16–9.17). BHSObjectApplicability is the project's explicit Adams-tower encoding of these hypotheses, not a quoted definition from BHS. The mathematical source has been checked; identification with the project's specified coefficient object, Adams towers, quotient maps and realization remains part of the existing formal comparison obligations. No Lean proof obligation is discharged by this citation. Primary references: [BHS](https://arxiv.org/pdf/1910.14116v3): Theorem A.1(3a), equivalently Theorem 9.19(3a); pages PDF pp. 70, 77 (also p. 36); arXiv:1910.14116v3, 10 May 2022; final prepublication version, not published pagination
% Blueprint note: Detailed source attribution: (Adapted from \cite{BHS}, Theorem A.1(3a), equivalently Theorem 9.19(3a).)
\end{theorem}

\begin{proposition}[The source-to-selected-model infinite-lift comparisons]
\label{sem:Literature:selected-lifts}
\level{coarse}
\uses{sem:Literature:infinite-nu-window,sem:Literature:infinite-finite-window,sem:AdamsSynthetic:finite-quotient-separation,sem:Literature:realization-coordinates}
(Adapted from \cite{BHS}, Thm. A.1 and Cor. A.19.)\par{}

For $X$ in the chosen classical object family, a classical label has a common infinite-cycle representative if and only if it is the first-quotient image of an untruncated class on $\nu X$. For a class of bidegree $(m,w)$ and $w-m\le s$, Adams filtration at least $s$ is equivalent to divisibility by $\lambda^{m+s-w}$ from weight $m+s$. A class on an object in the chosen synthetic family has zero realization if and only if some finite power of $\lambda$ annihilates it.
% Blueprint note: These are separately registered source-to-selected-model comparison obligations; they do not assert that every selected object satisfies the complete-source hypotheses.
\lean{KIP126.Literature.Standard.selected_permanent_lift,KIP126.Literature.Standard.selected_filtration_lambda,KIP126.Literature.Standard.selected_realization_kernel}
% Blueprint note: Source verification (2026-10-18): Partial source correspondence. BHS proves the source result for E-nilpotent complete spectra with a strongly convergent Adams spectral sequence. The selected family includes the ordinary sphere, its nu-cofiber, the detector, and shifts; membership alone does not establish BHS completeness. This entry is a pending source-to-selected-model comparison, not a direct application of BHS. The mathematical source has been checked; identification with the project's specified coefficient object, Adams towers, quotient maps and realization remains part of the existing formal comparison obligations. No Lean proof obligation is discharged by this citation. The realization-kernel clause is a consequence of the full synthetic category's localization telescope; its transport to this project's realization is a separate pending comparison. Primary references: [BHS](https://arxiv.org/pdf/1910.14116v3): Theorem A.1, permanent-cycle paragraph preceding (2a); proof of Theorem 9.19(2)–(4); pages PDF pp. 70, 76–77; arXiv:1910.14116v3, 10 May 2022; final prepublication version, not published pagination [BHS](https://arxiv.org/pdf/1910.14116v3): Corollary A.19; pages PDF p. 76; arXiv:1910.14116v3, 10 May 2022; final prepublication version, not published pagination [Pst](https://arxiv.org/pdf/1803.01804v3): Proposition 4.33; Proposition 4.40; Remark 4.41; pages PDF pp. 46–49; arXiv:1803.01804v3, 10 November 2022; not published pagination
% Blueprint note: Detailed source attribution: (Adapted from \cite{BHS}, Theorem A.1, the permanent-cycle paragraph, and Corollary A.19; \cite{Pst}, Remark 4.41. Source-to-selected-model comparison pending.)
\end{proposition}

\begin{proposition}[The selected-model realization comparisons]
\label{sem:Literature:selected-realization}
\level{coarse}
\uses{sem:Literature:selected-lifts,sem:Literature:realization-coordinates,sem:AdamsSynthetic:e2-comparison}
(Adapted from \cite{BHS}, Thm. A.1(2)--(3).)\par{}

Let $X$ belong to the chosen classical object family. A lift $a$ of a common permanent label $x$ surviving nontrivially to $E_{r+1}$ satisfies $\lambda^{r-1}a\ne0$ for $r\ge1$. If $x$ is a nonzero permanent label, it detects $\operatorname{Re}(a)$; each prescribed classical class detected by $x$ is the realization of such a lift. If a common infinite-cycle label $y$ is hit on page $r\ge2$, it has a lift annihilated by $\lambda^{r-1}$.
% Blueprint note: The producer must compare completed source objects with the exact selected objects. No blanket completeness assertion replaces those comparison responsibilities.
\lean{KIP126.Literature.Standard.selected_realization_lifetime,KIP126.Literature.Standard.selected_realization_detection,KIP126.Literature.Standard.selected_realization_prescribed_lift,KIP126.Literature.Standard.selected_realization_boundary_lift}
% Blueprint note: Source verification (2026-10-18): Partial source correspondence. BHS proves the source result for E-nilpotent complete spectra with a strongly convergent Adams spectral sequence. The selected family includes the ordinary sphere, its nu-cofiber, the detector, and shifts; membership alone does not establish BHS completeness. This entry is a pending source-to-selected-model comparison, not a direct application of BHS. The mathematical source has been checked; identification with the project's specified coefficient object, Adams towers, quotient maps and realization remains part of the existing formal comparison obligations. No Lean proof obligation is discharged by this citation. Primary references: [BHS](https://arxiv.org/pdf/1910.14116v3): Theorem A.1(2a), equivalently Theorem 9.19(2a); pages PDF pp. 70, 77 (also p. 36); arXiv:1910.14116v3, 10 May 2022; final prepublication version, not published pagination [BHS](https://arxiv.org/pdf/1910.14116v3): Theorem A.1(2b) and (3b), equivalently Theorem 9.19(2b),(3b); pages PDF pp. 70, 77 (also p. 36); arXiv:1910.14116v3, 10 May 2022; final prepublication version, not published pagination [BHS](https://arxiv.org/pdf/1910.14116v3): Theorem A.1(3a), equivalently Theorem 9.19(3a); pages PDF pp. 70, 77 (also p. 36); arXiv:1910.14116v3, 10 May 2022; final prepublication version, not published pagination
% Blueprint note: Detailed source attribution: (Adapted from \cite{BHS}, Theorem A.1(2a),(2b),(3a),(3b). Source-to-selected-model comparison pending.)
\end{proposition}

\section{The Kervaire survival criterion}

\begin{definition}[The distinguished synthetic theta-five for the criterion]
\label{sem:Literature:bx-choice}
\level{coarse}
\uses{sem:Literature:theta-five,sem:Literature:hopf-classes}
(Adapted from \cite{BurklundXu}, Prop. 7.19.)\par{}

Fix the distinguished class $\theta_5^{\mathrm{BX}}\in\pi_{62,64}S^{0,0}$ of order two and first-quotient label $h_5^2$. Let $\eta$ be the normalized synthetic Hopf class with label $h_1$. The products $\eta(\theta_5^{\mathrm{BX}})^2$ and $\lambda\eta(\theta_5^{\mathrm{BX}})^2$ use these choices.
% Blueprint note: Transport from the distinguished BX theta-five to arbitrary detected choices is a separate Main argument.
\lean{KIP126.Literature.Standard.bxTheta5,KIP126.Literature.Standard.bx_theta_label,KIP126.Literature.Standard.bx_order_two,KIP126.Literature.Standard.bx_eta_label,KIP126.Literature.Standard.etaThetaSquare,KIP126.Literature.Standard.lambdaEtaThetaSquare}
% Blueprint note: Source verification (2026-10-18): Source adaptation checked. This is the project-selected BX source choice, not a claim for every class with the same label. BX uses bidegree (stem, Adams filtration); the project uses (stem, weight), with weight = stem + filtration. Thus the source theta in degree (62,2) is displayed here in degree (62,64), and the source total boundary target is regraded to the project shift (1,-1). The finite statement retains r >= 1. Source-to-model identification and the registered Lean proof obligations remain separate from this citation check. Primary references: [BurklundXu](https://arxiv.org/pdf/2302.11869v3): Proposition 7.19 and its proof; Theorem 7.20 and its proof; pages 54; arXiv:2302.11869v3, 9 November 2024; accepted version [LWX126](https://arxiv.org/abs/2412.10879): Theorem 7.3, Remarks 7.4–7.5; Example 5.8(4); proof of Lemma 7.16; pages 42, 28, 49–50; Repository AimPaper/main.tex and AimPaper/2412.10879.pdf snapshot; no unstated version substitution
% Blueprint note: Detailed source attribution: (Adapted from \cite{BurklundXu}, Proposition 7.19 and the proof of Theorem 7.20; see \cite{LWX126}, Notation 7.2 and Remark 7.5.)
\end{definition}

\begin{theorem}[The finite and untruncated survival criteria]
\label{sem:Literature:bx-survival}
\level{coarse}
\uses{sem:Literature:bx-choice,sem:AdamsSynthetic:lambda-quotients,sem:AdamsSynthetic:standard-classes}
(Adapted from \cite{BurklundXu}, Prop. 7.19.)\par{}

For $r\ge1$, $h_6^2$ survives to $E_{r+3}$ if and only if the image of $\eta(\theta_5^{\mathrm{BX}})^2$ in the finite quotient of length $r$ vanishes. In the untruncated sphere, $h_6^2$ is a nonzero permanent cycle if and only if
\[\lambda\eta(\theta_5^{\mathrm{BX}})^2=0.\]
% Blueprint note: The finite statement uses the registered quotient-length-r convention for positive natural-number parameters. The full paper formulation and its indexing are distinct.
\lean{KIP126.Literature.Standard.bx_finite,KIP126.Literature.Standard.bx_untruncated}
% Blueprint note: Source verification (2026-10-18): Source adaptation checked. This is the project-selected BX source choice, not a claim for every class with the same label. BX uses bidegree (stem, Adams filtration); the project uses (stem, weight), with weight = stem + filtration. Thus the source theta in degree (62,2) is displayed here in degree (62,64), and the source total boundary target is regraded to the project shift (1,-1). The finite statement retains r >= 1. Source-to-model identification and the registered Lean proof obligations remain separate from this citation check. Primary references: [BurklundXu](https://arxiv.org/pdf/2302.11869v3): Proposition 7.19 and its proof; Theorem 7.20 and its proof; pages 54; arXiv:2302.11869v3, 9 November 2024; accepted version [LWX126](https://arxiv.org/abs/2412.10879): Theorem 7.3, Remarks 7.4–7.5; Example 5.8(4); proof of Lemma 7.16; pages 42, 28, 49–50; Repository AimPaper/main.tex and AimPaper/2412.10879.pdf snapshot; no unstated version substitution
% Blueprint note: Detailed source attribution: (Adapted from \cite{BurklundXu}, Proposition 7.19 and its proof; recalled in \cite{LWX126}, Theorem 7.3 and Remark 7.4.)
\end{theorem}

\begin{theorem}[The total boundary of the h-six square]
\label{sem:Literature:bx-total-boundary}
\level{coarse}
\uses{sem:Literature:bx-choice,sem:AdamsSynthetic:lambda-quotients,sem:AdamsSynthetic:standard-classes}
(Adapted from \cite{BurklundXu}, proof of Prop. 7.19.)\par{}

The total boundary of the first-quotient class labeled $h_6^2$ is
\[\delta(h_6^2)=\lambda\eta(\theta_5^{\mathrm{BX}})^2.\]
\lean{KIP126.Literature.Standard.deltaH6Square,KIP126.Literature.Standard.h6FirstQuotient,KIP126.Literature.Standard.h6FirstQuotient_label,KIP126.Literature.Standard.bx_total_boundary}
% Blueprint note: Source verification (2026-10-18): Source adaptation checked. This is the project-selected BX source choice, not a claim for every class with the same label. BX uses bidegree (stem, Adams filtration); the project uses (stem, weight), with weight = stem + filtration. Thus the source theta in degree (62,2) is displayed here in degree (62,64), and the source total boundary target is regraded to the project shift (1,-1). The finite statement retains r >= 1. Source-to-model identification and the registered Lean proof obligations remain separate from this citation check. Primary references: [BurklundXu](https://arxiv.org/pdf/2302.11869v3): Proposition 7.19 and its proof; Theorem 7.20 and its proof; pages 54; arXiv:2302.11869v3, 9 November 2024; accepted version [LWX126](https://arxiv.org/abs/2412.10879): Theorem 7.3, Remarks 7.4–7.5; Example 5.8(4); proof of Lemma 7.16; pages 42, 28, 49–50; Repository AimPaper/main.tex and AimPaper/2412.10879.pdf snapshot; no unstated version substitution
% Blueprint note: Detailed source attribution: (Adapted from \cite{BurklundXu}, proof of Proposition 7.19; regraded as in \cite{LWX126}, Remark 7.5.)
\end{theorem}

\section{May push-pull and secondary operations}

\begin{definition}[Exact tensor functors in the same synthetic category]
\label{sem:Literature:may-tensor}
\level{coarse}
\uses{sem:AdamsSynthetic:categories,frontier:AdamsSynthetic:d505f4e20090}
(Adapted from \cite{Pst}, Prop. 4.2; \cite{May01}, TC1--TC2.)\par{}

Left and right tensoring by a synthetic object commute with suspension and preserve distinguished triangles. Equip both functors with these shift isomorphisms and exactness structures; together they define the May tensor context.
\lean{KIP126.Literature.Standard.may_leftShift,KIP126.Literature.Standard.may_rightShift,KIP126.Literature.Standard.may_leftExact,KIP126.Literature.Standard.may_rightExact,KIP126.Literature.Standard.mayData}
% Blueprint note: Source verification (2026-10-18): Source adaptation checked. These choices realize integral suspension comparisons for the existing fixed synthetic tensor; TC1–TC2, not Lemma 4.6, are the relevant framework. Compatibility of this concrete choice remains to be constructed. Colimit preservation implies exactness for the synthetic stable category; the Lean declarations require this for the exact suspension structures selected above. Packages the selected left/right shift and exactness structures. The package contains TC1–TC2-style data only; it does not itself imply TC3. Primary references: [May01](https://www.math.uchicago.edu/~may/PAPERS/AddJan01.pdf): Axioms TC1–TC2; Remark 4.2; pages 11; Author-hosted January 2001 preprint, 31 pages; numbered statements retain published numbering [Pst](https://arxiv.org/pdf/1803.01804v3): Proposition 4.2; Remark 4.5; Definition 4.6; pages 38–39; arXiv:1803.01804v3, 10 November 2022
% Blueprint note: Detailed source attribution: (Adapted from \cite{Pst}, Proposition 4.2, and \cite{May01}, Axioms TC1--TC2 and Remark 4.2.)
\end{definition}

\begin{theorem}[May push-pull data and the signed boundary identity]
\label{sem:Literature:may-push-pull}
\level{coarse}
\uses{sem:Literature:may-tensor}
(Adapted from \cite{May01}, Lem. 4.6.)\par{}

For two cofiber triangles in the May tensor context, their tensor-product diagram admits May push-pull data and a signed boundary identity. If multiplication by two annihilates the relevant boundary image, the signed identity gives the corresponding unsigned identity.
\lean{KIP126.Literature.Standard.may_pushpull,KIP126.Literature.Standard.may_signed_boundary}
% Blueprint note: Source verification (2026-10-18): Source adaptation checked. Lean proves the signed boundary identity from may_pushpull. Removing the sign requires the image under the relevant group homomorphism to have exponent two, not just F2 Adams coordinates. The exact interface is the signed TC3 pushpull consequence for these selected tensor/suspension conventions, retaining the common vertex, all comparison squares, and represented homotopy lift. Primary references: [May01](https://www.math.uchicago.edu/~may/PAPERS/AddJan01.pdf): Axiom TC3; Lemma 4.6; pages 12–14; Author-hosted January 2001 preprint, 31 pages; numbered statements retain published numbering
% Blueprint note: Detailed source attribution: (Adapted from \cite{May01}, Axiom TC3 and Lemma 4.6; see \cite{LWX126}, Lemma 6.11.)
\end{theorem}

\begin{definition}[Toda witnesses and their indeterminacy]
\label{sem:Literature:toda-relations}
\level{coarse}
\uses{sem:AdamsSynthetic:categories,sem:AdamsSynthetic:shift-coherence}
(Adapted from \cite{ChristensenFrankland2017}, Def. 3.1 and Lem. 3.4.)\par{}

A triple Toda witness consists of a distinguished cofiber triangle, null composites, an extension over the cofiber, and a boundary equation giving its value. A witness implies composability. For composable data, witnesses exist, and their values form a coset of the full Toda indeterminacy subgroup. In the synthetic sphere, use the fixed shift identifications.
\lean{KIP126.Literature.Standard.TodaRelation,KIP126.Literature.Standard.TripleToda,KIP126.Literature.Standard.toda_composable,KIP126.Literature.Standard.toda_exists_relation,KIP126.Literature.Standard.toda_relation_iff_indeterminacy}
% Blueprint note: Source verification (2026-10-18): Source adaptation checked. The project uses the iterated-cofiber convention, with a fixed suspension identification for the synthetic TripleToda specialization. Existence uses exactness of a distinguished triangle. Full two-summand indeterminacy is retained; it is not replaced by a chosen bracket value. Primary references: [ChristensenFrankland2017](https://arxiv.org/pdf/1510.09216v3): Definition 3.1, Lemmas 3.4–3.5 and 5.12; pages 8, 11, 21; arXiv:1510.09216v3, 12 August 2017; closely matches AGT 17 (2017), 2687–2735 [LWX126](https://arxiv.org/abs/2412.10879): Theorem 7.3, Remarks 7.4–7.5; Example 5.8(4); proof of Lemma 7.16; pages 42, 28, 49–50; Repository AimPaper/main.tex and AimPaper/2412.10879.pdf snapshot; no unstated version substitution
% Blueprint note: Detailed source attribution: (Adapted from \cite{ChristensenFrankland2017}, Definition 3.1 and Lemma 3.4; used in \cite{LWX126}, proof of Lemma 7.16.)
\end{definition}

\begin{proposition}[Naturality, juggling, suspension, and tensoring for Toda relations]
\label{sem:Literature:toda-functoriality}
\level{coarse}
\uses{sem:Literature:toda-relations,sem:Literature:may-tensor}
(Adapted from \cite{ChristensenFrankland2017}, Lem. 3.5 and Lem. 5.12.)\par{}

Pre- and postcomposition carry Toda witnesses to the corresponding witnesses, and factors can be absorbed into the first and last entries. Juggling, shuffle, and suspension preserve the cofiber relations. A compatible functor sends a witness to a witness. Tensoring on either side does the same when equipped with exactness and shift compatibility.
\lean{KIP126.Literature.Standard.toda_precompose,KIP126.Literature.Standard.toda_postcompose,KIP126.Literature.Standard.toda_absorb_first,KIP126.Literature.Standard.toda_absorb_last,KIP126.Literature.Standard.toda_juggling,KIP126.Literature.Standard.toda_shuffle_iff,KIP126.Literature.Standard.toda_suspension_iff,KIP126.Literature.Standard.toda_map,KIP126.Literature.Standard.toda_tensor_left,KIP126.Literature.Standard.toda_tensor_right}
% Blueprint note: Source verification (2026-10-18): Source adaptation checked. Pre/postcomposition and absorption are containments, not equality of arbitrary brackets. Juggling/shuffle additionally assume the octahedral axiom and retain the shifted negative sign. Suspension retains (-1)^n and actual shift comparisons. Exact-functor and tensor claims require CommShift and IsTriangulated for the same functor; their proofs transport the complete cofiber witness and claim containment only. Primary references: [ChristensenFrankland2017](https://arxiv.org/pdf/1510.09216v3): Definition 3.1, Lemmas 3.4–3.5 and 5.12; pages 8, 11, 21; arXiv:1510.09216v3, 12 August 2017; closely matches AGT 17 (2017), 2687–2735 [LWX126](https://arxiv.org/abs/2412.10879): Theorem 7.3, Remarks 7.4–7.5; Example 5.8(4); proof of Lemma 7.16; pages 42, 28, 49–50; Repository AimPaper/main.tex and AimPaper/2412.10879.pdf snapshot; no unstated version substitution
% Blueprint note: Detailed source attribution: (Adapted from \cite{ChristensenFrankland2017}, Definition 3.1 and Lemmas 3.5 and 5.12; the exact-functor and tensor transport laws are derived here.)
\end{proposition}

\begin{proposition}[The synthetic lift of h-zero]
\label{sem:Literature:synthetic-h-zero}
\level{coarse}
\uses{sem:Literature:hopf-classes,sem:Literature:one-line-differentials,sem:Literature:first-quotient-products}
(Adapted from \cite{LWX126}, Ex. 5.8(4); \cite{BHS}, Prop. A.20.)\par{}

There is a class $[h_0]\in\pi_{0,1}S^{0,0}$ with first-quotient label $h_0$ such that $\lambda[h_0]=2$ and $[h_0]\eta=0$. Moreover,
\[[h_0]a=0\qquad\text{for all }a\in\pi_{2,3}S^{0,0}.\]
\lean{KIP126.Literature.Standard.syntheticH0,KIP126.Literature.Standard.toda_h0_label,KIP126.Literature.Standard.toda_lambda_h0,KIP126.Literature.Standard.toda_h0_eta,KIP126.Literature.Standard.toda_low_indeterminacy}
% Blueprint note: Source verification (2026-10-18): Source adaptation checked. The fixed h-zero choice is normalized by the first-quotient label and lambda[h-zero]=2, as in AimPaper Example 5.8(4). BHS Proposition A.20 gives a complete low synthetic ring presentation, with relation (0) h-zero eta=0 and relation (9) tau h-zero=2. In stem two, the presentation forces pi_(2,3) to be generated by tau eta^2 (the h-zero factors vanish), so h-zero annihilates every element there. This last assertion is a derived low-degree consequence, not a verbatim statement of AimPaper Lemma 7.16. The project must still identify its selected constants with these normalized low generators. Primary references: [LWX126](https://arxiv.org/abs/2412.10879): Theorem 7.3, Remarks 7.4–7.5; Example 5.8(4); proof of Lemma 7.16; pages 42, 28, 49–50; Repository AimPaper/main.tex and AimPaper/2412.10879.pdf snapshot; no unstated version substitution [BHS](https://arxiv.org/pdf/1910.14116v3): Remark 9.5; Theorem A.8; Corollary A.9(1); Notation A.10, Corollary A.11; Proposition A.20, relations (0), (9), and Remark A.22; pages 32, 72–73, 78–79; published Acta page mapping supplied in bhs_pages/proposals.json; arXiv:1910.14116v3, as independently retrieved by the BHS audit worker
% Blueprint note: Detailed source attribution: (Adapted from \cite{LWX126}, Example 5.8(4), and \cite{BHS}, Proposition A.20, relations (0) and (9).)
\end{proposition}

\begin{theorem}[The low synthetic Toda value]
\label{sem:Literature:low-toda-comparison}
\level{coarse}
\uses{sem:Literature:toda-relations,sem:Literature:synthetic-h-zero,sem:Literature:hopf-classes}
(Adapted from \cite{LWX126}, proof of Lem. 7.16.)\par{}

In the synthetic sphere,
\[\eta^2\in\langle[h_0],\eta,[h_0]\rangle.\]
\lean{KIP126.Literature.Standard.toda_eta_squared}
% Blueprint note: Source verification (2026-10-18): Source adaptation checked. The paper states equality of the low bracket with eta squared; the project asks only membership in its cofiber-based relation, with degree (2,4). This low identity is used within the proof and is not the statement of the conditional Lemma 7.16. Primary references: [LWX126](https://arxiv.org/abs/2412.10879): Theorem 7.3, Remarks 7.4–7.5; Example 5.8(4); proof of Lemma 7.16; pages 42, 28, 49–50; Repository AimPaper/main.tex and AimPaper/2412.10879.pdf snapshot; no unstated version substitution
% Blueprint note: Detailed source attribution: (Adapted from \cite{LWX126}, proof of Lemma 7.16, p. 49.)
\end{theorem}

\begin{theorem}[The symmetric Toda value for an order-two class]
\label{sem:Literature:symmetric-toda-comparison}
\level{coarse}
\uses{sem:Literature:toda-relations,sem:Literature:hopf-classes,sem:Literature:synthetic-h-zero}
(Adapted from \cite{Toda1962}, Cor. 3.7; \cite{BauesMuro2009}, Thm. 1.3.)\par{}

For every $\theta\in\pi_{62,64}S^{0,0}$ with $2\theta=0$,
\[\lambda^2\eta\theta\in\langle2,\theta,2\rangle.\]
\lean{KIP126.Literature.Standard.toda_symmetric}
% Blueprint note: Source verification (2026-10-18): Partial source correspondence. Toda Corollary 3.7 gives the uniform classical order-two symmetric-bracket law. Baues--Muro Definition 1.2, relation (T6), and Theorem 1.3 supply the connective-ring-spectrum version with the image of the stable Hopf element. To apply it in the synthetic category, one uses the canonical stable enrichment (or the corresponding square-zero mapping module) and identifies that Hopf element in pi_(1,0) with lambda^2 eta using BHS Proposition A.20. AimPaper proof of Lemma 7.16 uses the displayed synthetic formula for theta-five. The bound Lean statement is broader: every theta in pi_(62,64) killed by two, without a theta-five detection hypothesis. The cited uniform law explains this extension, but the canonical-enrichment and selected-Hopf-class comparison is an additional source-to-project obligation, not a quoted assertion of Lemma 7.16. Primary references: [Toda1962](https://webhomes.maths.ed.ac.uk/~v1ranick/papers/toda.pdf): Corollary 3.7; pages 31 (right-hand page of PDF page 18); Composition Methods in Homotopy Groups of Spheres, original book scan [BauesMuro2009](https://idus.us.es/server/api/core/bitstreams/57be0e00-a434-4786-b7e3-83a7eeed4e7a/content): Definition 1.2, relation (T6), and Theorem 1.3; pages 3 in the author manuscript; Author manuscript arXiv:math/0612450v1, 15 December 2006; published Comm. Algebra 37 (2009), 56–82 [LWX126](https://arxiv.org/abs/2412.10879): Theorem 7.3, Remarks 7.4–7.5; Example 5.8(4); proof of Lemma 7.16; pages 42, 28, 49–50; Repository AimPaper/main.tex and AimPaper/2412.10879.pdf snapshot; no unstated version substitution [BHS](https://arxiv.org/pdf/1910.14116v3): Remark 9.5; Theorem A.8; Corollary A.9(1); Notation A.10, Corollary A.11; Proposition A.20, relations (0), (9), and Remark A.22; pages 32, 72–73, 78–79; published Acta page mapping supplied in bhs_pages/proposals.json; arXiv:1910.14116v3, as independently retrieved by the BHS audit worker
% Blueprint note: Detailed source attribution: (Adapted from \cite{Toda1962}, Corollary 3.7, and \cite{BauesMuro2009}, Definition 1.2 (T6) and Theorem 1.3; used synthetically in \cite{LWX126}, proof of Lemma 7.16. The uniform synthetic comparison remains to be checked.)
\end{theorem}

\section{Multiplicative comparison inputs}

\begin{definition}[Algebra structures on the positive lambda quotients]
\label{sem:Literature:quotient-algebras}
\level{coarse}
\uses{sem:AdamsSynthetic:lambda-quotients,sem:AdamsSynthetic:graded-hom}
(Adapted from \cite{LWX126}, Rem. 3.4.)\par{}

For each $q>0$, equip $S^{0,0}/\lambda^q$ with a commutative algebra structure in the symmetric synthetic category whose unit is the quotient inclusion. Its multiplication and the sphere tensor comparison induce a bigraded product. Multiplication by a projected sphere class equals the sphere action.
\lean{KIP126.Literature.Standard.quotientSyntheticSymmetric,KIP126.Literature.Standard.quotientAlgebra,KIP126.Literature.Standard.quotientAlgebra_commutative,KIP126.Literature.Standard.quotientAlgebra_unit,KIP126.Literature.Standard.quotientAlgebraProduct,KIP126.Literature.Standard.quotientAlgebra_sphere_action}
% Blueprint note: Source verification (2026-10-18): Source adaptation checked. The all-positive-power E-infinity assertion is explicitly Remark 3.4 of AimPaper. BHSmot B.5 directly concerns the first special fiber and C.22 has cellular/completion qualifications; no invented numbered all-q theorem is claimed. Primary references: [LWX126](https://arxiv.org/abs/2412.10879): Remark 3.4; pages 15; Repository-frozen AimPaper/main.tex and AimPaper/2412.10879.pdf [BHSmot](https://arxiv.org/pdf/2010.10325v2): Definition B.3; Lemma B.5; Proposition C.2; Proposition C.22; pages 69–70, 75; arXiv:2010.10325v2, 30 June 2022 [Pst](https://arxiv.org/pdf/1803.01804v3): Corollary 4.30; Theorem 4.46; Lemma 4.56; pages 46, 50–51, 54; arXiv:1803.01804v3, 10 November 2022 [Pst](https://arxiv.org/pdf/1803.01804v3): Proposition 4.2; Remark 4.5; Definition 4.6; pages 38–39; arXiv:1803.01804v3, 10 November 2022 [Pst](https://arxiv.org/pdf/1803.01804v3): Definition 4.6; pages 39; arXiv:1803.01804v3, 10 November 2022
% Blueprint note: Detailed source attribution: (Adapted from \cite{LWX126}, Remark 3.4; see \cite{BHSmot}, Definition B.3, Lemma B.5, and Proposition C.22, and \cite{Pst}, Proposition 4.2.)
\end{definition}

\begin{proposition}[The first quotient identifies products with Ext products]
\label{sem:Literature:first-quotient-products}
\level{coarse}
\uses{sem:Literature:quotient-algebras,sem:AdamsSynthetic:e2-comparison,sem:AdamsSynthetic:first-quotient-shifts,frontier:AdamsSynthetic:f05463b82afe}
(Adapted from \cite{Pst}, Thm. 4.46 and Lem. 4.56.)\par{}

Under the first-quotient equivalence, multiplication of classes in bidegrees $(t-s,t)$ and $(t'-s',t')$ is the classical sphere $E_2$ product in bidegree $(s+s',t+t')$. The action of an untruncated sphere class on a first-quotient class agrees with the same product.
\lean{KIP126.Literature.Standard.sphereFirstQuotient,KIP126.Literature.Standard.quotientAlgebra_first_quotient,KIP126.Literature.Standard.first_quotient_action}
% Blueprint note: Source verification (2026-10-18): Source adaptation checked. Symmetric monoidality gives multiplicativity of the canonical Ext comparison, with s,t,s-prime,t-prime nonnegative and stem t-s. Identification with the project’s selected firstQuotient and sphereAlgebra product is still an explicit pending comparison. This packages the already selected sphere first-quotient equivalence; its external origin is the source Ext comparison, while this exact additive equivalence is project data. The Lean proof composes quotientAlgebra_sphere_action with quotientAlgebra_first_quotient. It is not a further external theorem or a completed upstream comparison. Primary references: [Pst](https://arxiv.org/pdf/1803.01804v3): Corollary 4.30; Theorem 4.46; Lemma 4.56; pages 46, 50–51, 54; arXiv:1803.01804v3, 10 November 2022
% Blueprint note: Detailed source attribution: (Adapted from \cite{Pst}, Theorem 4.46 and Lemma 4.56, using the first-quotient algebra of Corollary 4.30.)
\end{proposition}

\begin{proposition}[Filtration and detection of actual products]
\label{sem:Literature:filtered-product-detection}
\level{coarse}
\uses{frontier:AdamsSynthetic:791f3c0e6a4a,sem:AdamsSynthetic:selected-convergence,sem:Literature:finite-quotient-labels}
(Adapted from \cite{Moss}, §1; \cite{BHS}, Rem. 9.5.)\par{}

A sphere class of filtration at least $s$ acting on a class of filtration at least $s'$ gives a class of filtration at least $s+s'$. On each positive finite quotient, if the two classes are detected by specified $E_2$ labels, their action is detected by the product of those labels, with the corresponding degree identifications. The classical product of detected sphere classes is likewise detected by the product of their classical $E_2$ labels.
\lean{KIP126.Literature.Standard.action_filtration,KIP126.Literature.Standard.finite_action,KIP126.Literature.Standard.classical_product_detection}
% Blueprint note: Source verification (2026-10-18): Partial source correspondence. The three bindings are general comparison statements for the project-selected Adams towers, not a single numbered theorem in AimPaper. Moss Section 1 supplies composition pairings in classical Adams systems, while BHS Remark 9.5 and the canonical towers in Appendix A explain the synthetic monoidal setting. The exact compatibility of the selected finite-quotient E2 coordinates, action and abutment for every positive quotient length, and the action-filtration statement on all shifts and iterated quotients, must still be supplied. No upper bound is inserted into the frozen formal range; this citation check does not silently narrow it to the paper examples. Primary references: [Moss](https://gdz.sub.uni-goettingen.de/download/pdf/PPN266833020_0115/LOG_0047.pdf): Section 1; Theorem (1.2) and conditions (1.3); pages 283–285; Original Math. Z. 115 (1970), 283–310; GDZ digitized journal scan LOG_0047 [BHS](https://arxiv.org/pdf/1910.14116v3): Remark 9.5; Theorem A.8; Corollary A.9(1); Notation A.10, Corollary A.11; Proposition A.20, relations (0), (9), and Remark A.22; pages 32, 72–73, 78–79; published Acta page mapping supplied in bhs_pages/proposals.json; arXiv:1910.14116v3, as independently retrieved by the BHS audit worker [LWX126](https://arxiv.org/abs/2412.10879): Theorem 7.3, Remarks 7.4–7.5; Example 5.8(4); proof of Lemma 7.16; pages 42, 28, 49–50; Repository AimPaper/main.tex and AimPaper/2412.10879.pdf snapshot; no unstated version substitution
% Blueprint note: Detailed source attribution: (Adapted from the multiplicative framework of \cite{Moss}, Section 1, and \cite{BHS}, Remark 9.5 and Appendix A; the selected filtration and finite-quotient comparisons remain explicit obligations.)
\end{proposition}


\section{Classical inputs for the Adams second line}
\begin{definition}[Products of standard one-line classes]
\label{sem:Literature:two-line-products}
\level{coarse}
(From \cite{Adams60}, Theorem 2.5.1(2), p.~45.)\par{}

Write $H_{i,j}=h_i h_j\in E_2^{2,2^i+2^j}(S^0)$ for the actual bilinear product of the fixed standard classes. Its degree is the sum of the two one-line bidegrees; no independent label is chosen.
\lean{KIP126.Literature.Standard.hiProduct}
\end{definition}

\begin{theorem}[Classical permanence and exclusions on the second line]
\label{sem:Literature:two-line-classical-inputs}
\level{coarse}
The following classical source results are adapted to the fixed $2$-primary sphere Adams spectral sequence.
\begin{enumerate}
\item (Adapted from \cite{MayThesis1964}, Appendix A; Theorem II.7.9.)\par{}

The three products $H_{0,2},H_{0,3},H_{2,4}$ survive nontrivially to $E_\infty$.
\item (Adapted from \cite{MayThesis1964}, Appendix A; Theorem II.7.9.)\par{}

For $0\leq j\leq3$, the square $h_j^2$ survives nontrivially.
\item (Adapted from \cite{MahowaldTangora}, Section 8; \cite{BMMS1986}, VI, Prop. 1.18(ii).)\par{}

The square $h_4^2$ survives nontrivially.
\item (Adapted from \cite{Mahowaldetaj}, Theorem 1', p.~249.)\par{}

For every $j\geq3$, $H_{1,j}$ survives nontrivially.
\item (Adapted from \cite{IWX}, Lemma 5.10(1), p.~143.)\par{}

There is a nonzero $d_3$ from $H_{2,5}$ at $(2,36)$ to degree $(5,38)$.
\item (Adapted from \cite{IWX}, Table 7, p.~225.)\par{}

There is a nonzero $d_4$ from $H_{3,5}$ at $(2,40)$ to degree $(6,43)$.
\item (Adapted from \cite{IWX}, Table 6, p.~223.)\par{}

There is a nonzero $d_3$ from $H_{3,6}$ at $(2,72)$ to degree $(5,74)$.
\item (Adapted from \cite{HHR}, Theorem 1.1, p.~2.)\par{}

For every $j\geq7$, $h_j^2$ does not survive nontrivially to $E_\infty$.
\item (Adapted from \cite{IWX}, Table 6, p.~222.)\par{}

There is a nonzero $d_3$ from $H_{0,4}$ at $(2,17)$ to degree $(5,19)$.
\end{enumerate}
\lean{KIP126.Literature.Standard.may_two_line_products_permanent,KIP126.Literature.Standard.may_low_squares_permanent,KIP126.Literature.Standard.mahowald_tangora_h4_square_permanent,KIP126.Literature.Standard.mahowald_eta_family_permanent,KIP126.Literature.Standard.mahowald_tangora_h2_h5_d3,KIP126.Literature.Standard.mahowald_tangora_h3_h5_d4,KIP126.Literature.Standard.iwx_h3_h6_d3,KIP126.Literature.Standard.hhr_high_squares_not_permanent,KIP126.Literature.Standard.iwx_h0_h4_d3}
% Blueprint note: Clauses 1--2 detailed source attribution: \cite{MayThesis1964}, Appendix A, pp.~A.1--A.2, and Theorem II.7.9, pp.~II-7.6--II-7.7, checked directly in the original 1964 Princeton thesis. The chart on A.2 contains h0^2, h1^2, h0 h2, h2^2, h0 h3 and h3^2; its source/target gaps and the standard connectivity argument exclude differentials. Theorem II.7.9(13) identifies nu-star with the E-infinity survivor h2 h4; parts (2), (3), (9) identify the nonzero Hopf-map squares. The h0^2 case follows from the chart and grading, not a direct assertion of that theorem. Remarks II.7.2, p.~II-7.2, identifies multiplication by h0 on E-infinity with multiplication by two. The existing AimPaper key Maythesis denotes the distinct 1966 Journal of Algebra article; it is not the thesis containing these locators. Original scan: [May thesis](https://www.math.ens.psl.eu/~obenoist/textes/May.pdf).
% Blueprint note: Clause 3 detailed source attribution: \cite{MahowaldTangora}, Section 8, is the historical source; this locator is explicitly corroborated in \cite{Browder}, p.~178, immediately after Theorem 7.1. The original MT scan has not been retrieved, so no independently verified MT theorem number is asserted. A directly inspected independent original proof is \cite{BMMS1986}, Chapter VI, Proposition 1.18(ii), p.~183: [Bruner's chapter](https://www.sas.rochester.edu/mth/sites/doug-ravenel/otherpapers/BMMS7.pdf). The public declaration's historical name is retained; its fixed-model proof is pending.
% Blueprint note: Clauses 5--6 detailed source attribution: The directly checked computations are \cite{IWX}, Lemma 5.10(1), p.~143, and Table 6, p.~222, row (34,2,18), giving motivic d3(h2 h5)=tau h1 d1; and Table 7, p.~225, row (38,2,20), giving d4(h3 h5)=h0 x. The classical passage uses Corollary 1.2 and its proof, p.~111. The authors' original classical data \cite{IWXCharts}, Adams-classical-E3.csv, physical rows 115 and 135 including the header, explicitly records the same d3 and d4 with nonzero targets h1 d1 and h0 x (rows 111 and 130). Historical sources are \cite{MahowaldTangora} and \cite{BarrattMahowaldTangora}; the original 1967/1970 scans have not been independently inspected, and historical attribution of the h2 h5 differential between them is not resolved here. The primary IWX recomputations, rather than an AimPaper report, supply the checked mathematical source. This does not silently rename the frozen declarations or discharge their source-to-fixed-model obligations.
% Blueprint note: Clause 7 detailed source attribution: \cite{IWX}, Theorem 5.7, p.~142, and Table 6, p.~223, row (70,2,36), give motivic d3(h3 h6)=h0 p-prime. Classical comparison is Corollary 1.2 and its proof, p.~111. Nonzero classical source and target are independently explicit in the authors' fixed original chart release \cite{IWXCharts}, Adams-classical-E3.csv, physical rows 472 and 457 including the header. Its SHA-256 is 1ebfd24afd3da79a8d2be0ea05ca96f90a6b1b3c5b8edf429ef49eca07f3b37b. This avoids inferring a nonzero classical differential from a motivic table without checking tau-torsion. The Lean interface still needs the comparison with the project's selected classes and pages.
% Blueprint note: Clause 9 detailed source attribution: \cite{IWX}, Table 6, p.~222, first row (15,2,8), gives motivic d3(h0 h4)=h0 d0. Corollary 1.2, p.~111, and the authors' original classical data \cite{IWXCharts}, Adams-classical-E3.csv, physical rows 29 and 27 including the header, directly confirm the nonzero classical source and target. This is the separately authorized additional input needed to correct the omitted E3 candidate h0 h4; it does not assert that this class is a d2 boundary.
% Blueprint note: Source scope (2026-10-18): All nine statements concern the actual fixed standard classes and actual spectral sequence. Nonzero survival requires a common permanent representative and nonzero E-infinity image. The four differentials quantify over actual E2 target labels and require nonzero page targets; no arbitrary label is substituted. HHR Theorem 1.1 does not specify a differential length or target. Mahowald Theorem 1', p.~249, including the nonboundary argument following Theorem 2, and HHR Theorem 1.1, p.~2, were checked directly. The primary-source references and qualifications above replace indirect AimPaper-only attribution, but all nine source-to-fixed-model Lean proofs remain separately registered upstream obligations. The original 57 AimPaper statement blocks and bibliography are unchanged.
\end{theorem}

\section{Framed manifolds and the Kervaire invariant}

\begin{definition}[Closed smooth manifolds and stable framings]
\label{sem:Literature:closed-framed-manifolds}
\level{coarse}
\textit{(Geometric conventions of \cite{LWX126}, Introduction, and \cite{Browder}, Sections 1--3.)}
A closed smooth $n$-manifold is a compact Hausdorff second-countable smooth
manifold without boundary, modelled on $\mathbb R^n$. A stable tangent
framing of such a manifold $M$ consists of a nonnegative integer $r$ and a
smooth global frame of $TM\oplus\varepsilon^r$: at each point the $n+r$
sections form a basis. A closed stably framed manifold is a pair of a
closed smooth manifold and a specified stable tangent framing.
No connectedness or unstabilized parallelizability is imposed.
\lean{KIP126.Literature.Geometry.ClosedSmoothManifold,KIP126.Literature.Geometry.StableFraming,KIP126.Literature.Geometry.ClosedStablyFramedManifold}
% Blueprint note: These are actual mathlib manifold and smooth vector-bundle data, with no construction holes. Their representations are module-private: ordinary imports expose the three geometric object types without changing the analytic instance environment of the older modules. A full-body representation audit verifies the actual compactness, smoothness and frame data and their independence from project spectral objects. The stable tangent/normal comparison needed by the literature theorem belongs to the pending Browder input.
\end{definition}

\begin{definition}[The framing-dependent Kervaire invariant]
\label{sem:Literature:kervaire-invariant}
\level{coarse}
\uses{sem:Literature:closed-framed-manifolds}
\textit{(From \cite{Browder}, Section 3, especially Theorem 3.2, pp.~166--168; see also \cite{LWX126}, Introduction.)}
For a closed smooth stably framed $(4k+2)$-manifold $(M,F)$, write
$\operatorname{KI}(M,F)\in\mathbb F_2$ for the Arf invariant of the
quadratic refinement of the middle-dimensional mod-two intersection pairing
induced by the specified framing. The invariant depends on the framed
manifold, not just its dimension.
\lean{KIP126.Literature.Geometry.kervaireInvariant}
% Blueprint note: Pending construction, explicitly authorized on 2026-10-18: construct the middle-dimensional pairing, its framing-induced quadratic refinement and its Arf invariant on the actual manifold model. This function remains a registered construction hole. Its type fixes the manifold and framing inputs and mod-two output, but the intended geometric construction is not yet implemented. It must not be filled by a dimension test or an Adams-survival predicate.
\end{definition}

\begin{definition}[Existence of a framed manifold of Kervaire invariant one]
\label{sem:Literature:kervaire-existence}
\level{coarse}
\uses{sem:Literature:kervaire-invariant}
For a natural number $n$, let $G(n)$ mean that there is a natural number
$k$ with $n=4k+2$ and a closed smooth stably framed $(4k+2)$-manifold
$(M,F)$ satisfying $\operatorname{KI}(M,F)=1$.
\lean{KIP126.Literature.Geometry.KervaireInvariantOneExists}
% Blueprint note: This is an existential definition over actual geometric data; it does not define existence by Browder's criterion. Its local definition is complete, but it depends on the pending invariant construction.
\end{definition}
